\documentclass[11pt,a4paper,oneside]{report}
\usepackage[T1]{fontenc}
\usepackage[utf8]{inputenc}
\usepackage{amsmath,amsthm,mathtools}
\usepackage{newtxtext,newtxmath}
\usepackage[margin=27mm,headheight=15pt,headsep=7mm,footskip=12mm]{geometry}
\usepackage{microtype}
\usepackage{booktabs,longtable,array,tabularx}
\usepackage{xcolor,listings,enumitem,needspace}
\usepackage{fancyhdr,titlesec,etoolbox}
\usepackage[unicode,pdfencoding=auto,psdextra]{hyperref}
\usepackage[nameinlink,noabbrev]{cleveref}
\definecolor{ink}{HTML}{162D43}
\definecolor{muted}{HTML}{53616D}
\definecolor{papergray}{HTML}{F5F6F7}
\hypersetup{colorlinks=true,linkcolor=ink,citecolor=ink,urlcolor=ink,
 pdftitle={Quantified CUE: A Set-Theoretic Constraint Calculus},
 pdfauthor={Aram Hăvărneanu},
 pdfsubject={Working technical proposal: quantification, function types, dependency, and abstraction},
 pdfkeywords={CUE, semantic subtyping, hyperdoctrine, polymorphism, existential types}}
\urlstyle{same}
\setlength{\parindent}{0pt}
\setlength{\parskip}{3pt plus 1pt minus 1pt}
\AtBeginDocument{
 \setlength{\abovedisplayskip}{7pt plus 2pt minus 2pt}
 \setlength{\belowdisplayskip}{7pt plus 2pt minus 2pt}
 \setlength{\abovedisplayshortskip}{3pt plus 1pt}
 \setlength{\belowdisplayshortskip}{5pt plus 1pt minus 1pt}
}
\setlength{\emergencystretch}{2em}
\setcounter{secnumdepth}{2}
\setcounter{tocdepth}{1}
\setlist{nosep,leftmargin=1.6em}
\pagestyle{fancy}
\fancyhf{}
\fancyhead[L]{\small\textsc{Quantified CUE}}
\fancyhead[R]{\small\nouppercase{\leftmark}}
\fancyfoot[C]{\thepage}
\renewcommand{\headrulewidth}{0.3pt}
\renewcommand{\chaptermark}[1]{\markboth{\thechapter\quad #1}{}}
\titleformat{\chapter}[display]{\normalfont\Large\bfseries\color{ink}}{\chaptertitlename\ \thechapter}{5pt}{\LARGE}
\titlespacing*{\chapter}{0pt}{0pt}{17pt}
\titleformat{\section}{\large\bfseries\color{ink}}{\thesection}{0.6em}{}
\titleformat{\subsection}{\normalsize\bfseries}{\thesubsection}{0.6em}{}
\titlespacing*{\section}{0pt}{12pt}{5pt}
\titlespacing*{\subsection}{0pt}{9pt}{3pt}
\pretocmd{\section}{\Needspace{9\baselineskip}}{}{}
\pretocmd{\subsection}{\Needspace{7\baselineskip}}{}{}
\lstdefinelanguage{CUE}{
 morekeywords={forall,exists,in,func,let,for,if,package,import,seal,with,open,as,extern,match,case,return},
 sensitive=true,morecomment=[l]{//},morecomment=[s]{/*}{*/},morestring=[b]",
 alsoletter={\#}}
\lstset{language=CUE,basicstyle=\ttfamily\fontsize{8.5}{10.4}\selectfont,
 keywordstyle=\bfseries\color{ink},commentstyle=\color{muted},stringstyle=\color{ink},
 columns=fullflexible,keepspaces=true,showstringspaces=false,tabsize=4,
 breaklines=true,breakatwhitespace=false,upquote=true,
 backgroundcolor=\color{papergray},frame=single,rulecolor=\color{black!13},
 framerule=0.3pt,framesep=5pt,xleftmargin=5pt,xrightmargin=5pt,
 aboveskip=6pt,belowskip=7pt}
\newcommand{\code}[1]{\texttt{\detokenize{#1}}}
\newcommand{\sem}[1]{\llbracket #1\rrbracket}
\newcommand{\Pow}{\mathcal P}
\newcommand{\Env}{\operatorname{Env}}
\newcommand{\FV}{\operatorname{FV}}
\newcommand{\App}{\operatorname{App}}
\newcommand{\Run}{\operatorname{run}}
\newcommand{\Ok}{\mathsf{Ok}}
\newcommand{\Reject}{\mathsf{Reject}}
\newcommand{\Raise}{\mathsf{Raise}}
\newcommand{\Type}{\mathsf{Type}}
\newcommand{\Int}{\mathbb Z}
\newcommand{\Num}{\mathsf{Number}}
\newcommand{\Str}{\mathsf{String}}
\newcommand{\Bool}{\mathsf{Bool}}
\newcommand{\List}{\mathsf{List}}
\newcommand{\Nat}{\mathbb N}
\newcommand{\U}{\mathcal U}
\DeclareMathAlphabet{\profilecal}{OMS}{cmsy}{m}{n}
\newcommand{\Uone}{\ensuremath{\profilecal{U}_{1}}}
\newcommand{\Sone}{\ensuremath{\profilecal{S}_{1}}}
\newcommand{\Shigher}{\ensuremath{\profilecal{S}_{H}}}
\newcommand{\Dep}{\ensuremath{\profilecal{D}}}
\newcommand{\Abs}{\ensuremath{\profilecal{A}}}
\pdfstringdefDisableCommands{%
 \def\Uone{U1}%
 \def\Sone{S1}%
 \def\Shigher{SH}%
 \def\Dep{D}%
 \def\Abs{A}%
}
\newcommand{\Rel}{\mathsf{Rel}}
\newcommand{\todoob}{\textbf{Implementation obligation.}\ }
\newtheorem{definition}{Definition}[chapter]
\newtheorem{proposition}[definition]{Proposition}
\newtheorem{theorem}[definition]{Theorem}
\newtheorem{lemma}[definition]{Lemma}
\theoremstyle{remark}
\newtheorem{remark}[definition]{Remark}
\newcounter{example}[chapter]
\renewcommand{\theexample}{\thechapter.\arabic{example}}
\newcommand{\example}[2]{\refstepcounter{example}\Needspace{5\baselineskip}\par\textbf{Example \theexample\ --- #1.}\ \textit{#2}\par\nopagebreak[4]}
\newcommand{\decision}[1]{\par\textbf{Design rule.} #1\par}
\newcolumntype{Y}{>{\raggedright\arraybackslash}X}
\newcolumntype{P}[1]{>{\raggedright\arraybackslash}p{#1}}
\begin{document}
\hypersetup{pageanchor=false}
\begin{titlepage}
\vspace*{16mm}
{\small\textsc{Working technical proposal}\par}
\vspace{12mm}
{\fontsize{31}{35}\selectfont\bfseries\color{ink}Quantified CUE\par}
\vspace{5mm}
{\Large A Set-Theoretic Constraint Calculus\par}
\vspace{3mm}
{\large Polymorphic functions, dependent specifications,\par
and representation-independent modules\par}
\vspace{17mm}
{\large Aram H\u{a}v\u{a}rneanu\par}
\vspace{2mm}
{22 September 2026\par}
\vfill
\begin{minipage}{\textwidth}
\textbf{Abstract.}
CUE descriptions are interpreted as predicates on joint environments of possible
values. Existential projection explains instance formation and hidden witnesses;
universal projection adds obligations over arbitrary values and semantic types.
Both are adjoints to substitution, and both preserve refinement. Function types
are universal predicates on an operational application relation. Their variance,
intersection laws, and polymorphic instances follow from that interpretation.
The proposal defines quantified declarations, legal scope extrusion, a
conservative data-language embedding, and three language profiles. A worked
catalogue covers calling conventions, higher-order programming, dependent
invariants, and abstract modules. Existential modules use sealing and
equality-respecting simulations to support representation-independent evolution.
The reference model establishes refinement, function soundness, and
representation-independent replacement under explicit hypotheses. A residual
normalizer integrates type obligations with refinement-stable propagation and
demanded concrete validation.
\end{minipage}\par
\vspace{8mm}
{\small Comparison baseline: CUE proposal 4484 and its theory addendum,\par
repository revision \texttt{a949162aedb8}, inspected 21 September 2026.\par}
\end{titlepage}
\pagenumbering{roman}
\hypersetup{pageanchor=true}
\tableofcontents
\clearpage
\chapter*{Conventions}
\addcontentsline{toc}{chapter}{Conventions}
Equations and proofs use the semantic reference calculus. Code listings use
the proposed surface language and carry profile or extension labels. An expected
contradiction is written \code{_|_}; an unresolved proof or evaluation obligation
is written \emph{incomplete}.

\begin{tabularx}{\textwidth}{@{}lY@{}}
\toprule
Notation & Meaning\\\midrule
$V$ & A fixed set of operational inhabitants, including data and closures.\\
$\rho\in\Env(\Gamma)$ & An assignment to the free variables of context $\Gamma$.\\
$\sem{T}_\rho\subseteq V$ & The inhabitants of description $T$ at $\rho$.\\
$P\le Q$ & Inclusion of predicates, hence refinement toward more information.\\
$\top,\bot$ & The full subject set and the empty predicate.\\
$\U_n$ & A set of admissible semantic types at universe level $n$.\\
\Uone{} & Universal-only, rank-1 surface profile.\\
\Sone{} & Symmetric rank-1 profile, with scoped existential packages.\\
\Shigher{} & Full symmetric, predicative higher-rank profile.\\
\Dep{} & The dependent-value extension of the indicated profile.\\
\Abs{} & The opaque-module abstraction discipline.\\
\bottomrule
\end{tabularx}

Universal-only refers to the \emph{new explicit type-binding surface}. Ordinary
CUE instance variables retain their existential completion semantics in all three
profiles. Rank restrictions concern type quantifiers; dependence on ordinary
values is an independent design axis.

The forms \code{forall (A, B) T} and \code{exists A { ... }} bind variables
in a description. A quantified field \code{f(A, B): T} constrains one subject
at every instance. A parametric alias \code{S(A) = T} abbreviates a description
by capture-avoiding substitution; \code{S(U)} applies the alias. Instantiation
selection \code{f[U]} refers to the same subject at a universal instance.

The comparison uses proposal 4484 and its addendum at the cited repository
revision \cite{cueproposal,cueaddendum}. Stable CUE expressions retain their
existing elaboration and observation rules.

\clearpage
\pagenumbering{arabic}
\chapter{The existing constraint semantics}
\label{ch:base}
\section{Descriptions and completions}
CUE orders descriptions by subsumption and defines unification as their greatest
lower bound. At the level of successful data completions, the corresponding
set interpretation is inclusion, intersection, and union
\cite{cuespec}.

Fix a set $D$ of complete data inhabitants. A primitive description denotes a
subset of $D$; a concrete atom denotes a singleton. Write
\[
 \sem{1}=\{1\},\qquad \sem{\code{int}}=\Int,\qquad
 \sem{S\mathbin{\&}T}=\sem S\cap\sem T,\qquad
 \sem{S\mathbin{|}T}=\sem S\cup\sem T.
\]
An open record denotes the records satisfying its field predicates; a closed
record additionally constrains its permitted labels. Field absence is a
separate possibility from a present field with an unconstrained value.

An open declaration is more accurately a relation than a product of its field
marginals. For example:
\par\noindent\begin{minipage}{\linewidth}
\begin{lstlisting}
config: {
    x: int
    y: x + 1
}
\end{lstlisting}
\end{minipage}\par
has the relational interpretation
\[
 R(x,y)\equiv x\in\Int\ \land\ y=x+1.
\]
Its two fields cannot be interpreted as independent choices of integers.
Projection onto an exported interface hides the other variables existentially.
A complete exported record retains its exposed fields as the subject being
constrained.

\section{The sense in which variability is existential}
A field declaration \code{x: int} leaves a witness to be supplied by a completion.
Adding \code{x: >=0} removes negative witnesses. It imposes no obligation to
produce an implementation working uniformly for every integer.

More generally, if $\Delta$ names local witnesses and $\Gamma$ the exposed
interface, the description of that interface is
\[
 \sem{\exists\Delta\,P}
 =\{\rho\in\Env(\Gamma)\mid
       \exists\delta\in\Env(\Delta).\;P(\rho,\delta)\}.
\]
This is the existential reading of ordinary CUE variability. Built-in list and
pattern constraints may themselves validate all matching elements; a general
binder for an arbitrary semantic type is an additional facility.

Sharing precedes projection. Two contributions to one field elaborate to
$\exists x\,(P(x)\land Q(x))$. Two independently introduced witnesses elaborate
to $(\exists x\,P(x))\land(\exists y\,Q(y))$. These formulas carry different
sharing information.

\example{Independent instantiations of a description}{Existing CUE}
\par\noindent\begin{minipage}{\linewidth}
\begin{lstlisting}
#Cell: {value: int}
left:  #Cell & {value: 1}
right: #Cell & {value: 2}
\end{lstlisting}
\end{minipage}\par
Each instantiation has its own field cells. Within each cell, repeated
constraints share that cell. CUE references copy the associated expression and
rebind internal references to the corresponding copies \cite{cuerefs}.
The reference calculus therefore receives an elaborated graph with explicit
sharing, rather than equating every textual identifier with one global logic
variable.

\section{Universal implementation obligations}
A generic transformation needs one implementation whose behavior is valid for
arbitrary types. Write $\operatorname{Map}(A,B,m)$ for the requirement that $m$
applies a supplied transformation from $A$ to $B$ elementwise to a finite
$A$-list and returns the corresponding $B$-list. The two statements
\[
 \exists A,B,m.\;\operatorname{Map}(A,B,m),\qquad
 \exists m.\;\forall A,B.\;\operatorname{Map}(A,B,m)
\]
are different specifications. The first asks for some usable instance. The
second fixes an implementation and requires every instance. In particular,
existentially choosing $A=\varnothing$ can trivialize an input obligation.

This distinction applies to records as well as functions: a reusable module may
contain one operation that works at every element type, and a hidden
representation shared by several operations. The language needs to retain both
kinds of binding and their dependency order.

\section{The semantic layer of existing features}
The reference model is the \emph{completion layer}. The implementation also
retains closure provenance, scope graphs, field-presence information, and
preference information. A default-bearing description has successful
completions and a preferred subset. Default selection is an observation of that
additional structure. Comprehensions elaborate constraints using the host
language's established evaluation rules \cite{cuedefaults}.

Legacy programs retain their data observations and evaluation rules.
Refinement theorems apply to a fixed elaborated predicate and a fixed quantifier
dependency structure.

\chapter{Quantification as a set-theoretic operation}
\label{ch:quantifiers}
\section{Contexts and subjects}
\begin{definition}[Predicate fiber]
For a context $\Gamma$ with set of assignments $\Env(\Gamma)$, let
$\mathcal P_\Gamma=\Pow(\Env(\Gamma))$, ordered by inclusion. A type expression
with subject $z$ denotes a predicate in $\mathcal P_{\Gamma,z:V}$.
\end{definition}
The subject is the value described by an expression. A quantifier binds its
listed variables and leaves that subject fixed. A record field supplies a
projection of the enclosing record as the subject of its field expression.

A substitution $h:\Env(\Delta)\to\Env(\Gamma)$ induces reindexing
$h^*P=h^{-1}(P)$. It preserves all set unions and intersections.
For a fixed sort $S$, let
$\pi:\Env(\Gamma)\times S\to\Env(\Gamma)$ be the projection.
Define
\begin{align}
 \exists_\pi P
 &=\{\rho\mid\exists s\in S.\;(\rho,s)\in P\},\label{eq:exists}\\
 \forall_\pi P
 &=\{\rho\mid\forall s\in S.\;(\rho,s)\in P\}.\label{eq:forall}
\end{align}
These are predicates on the original interface, not families left outside the
semantic domain.

\begin{proposition}[Adjunctions]\label{prop:adjoints}
For predicates $P$ on the extended context and $Q$ on the base context,
\[
 \exists_\pi P\subseteq Q\ \Longleftrightarrow\ P\subseteq\pi^*Q,
 \qquad
 \pi^*Q\subseteq P\ \Longleftrightarrow\ Q\subseteq\forall_\pi P.
\]
\end{proposition}
\begin{proof}
The first equivalence expands to: every $(\rho,s)\in P$ has $\rho\in Q$.
The second expands to: for every $\rho\in Q$ and every $s\in S$,
$(\rho,s)\in P$.
\end{proof}
Thus $\exists_\pi\dashv\pi^*\dashv\forall_\pi$. This is the subset
hyperdoctrine underlying the proposal \cite{lawvere,mellies}.

\section{Semantic types and universes}
Fix a set $V$ of operational values. A universe $\U_n\subseteq\Pow(V)$ is a
set of admissible type interpretations. A valuation assigns
$\rho(A)\in\U_n$ to a type variable of sort $\Type(n)$.
Type operations are interpreted pointwise in that valuation:
\[
 \sem{A}_\rho=\rho(A),\quad
 \sem{T\land U}_\rho=\sem T_\rho\cap\sem U_\rho,\quad
 \sem{T\lor U}_\rho=\sem T_\rho\cup\sem U_\rho.
\]
A fixed-level admissible family contains the base types and the operations used
by the selected fragment. Universes are cumulative and closed at higher levels
under the permitted quantified constructions. Taking sufficiently large
set-sized families of subsets gives a reference interpretation; an implementation
uses finitely represented descriptions of selected members.

For a binder ranging over $\U_n$:
\begin{align}
 \sem{\forall A\,T}_\rho
 &=\bigcap_{S\in\U_n}\sem T_{\rho[A:=S]},\label{eq:foralltype}\\
 \sem{\exists A\,T}_\rho
 &=\bigcup_{S\in\U_n}\sem T_{\rho[A:=S]}.\label{eq:existstype}
\end{align}
The quantifier does not replace the subject by a type-indexed family of subjects.
These equations are a Curry-style, erased interpretation of quantified
\emph{constraints on one value}.

Stratification gives a predicative source discipline. Operational functions will
be members of a designated set of closures, rather than the full set-theoretic
exponential $B^A$. This is a Henkin-style choice of realizers and admissible
predicates. It lies outside the standard full-function-space extension obstructed
by Reynolds's theorem \cite{reynolds84}. The equations alone do not assert
relational parametricity; that property additionally concerns the operations
available to an abstractly typed client \cite{reynolds83}.

\section{Bounds and dependent ranges}
A type bound $A:T$ abbreviates the range
\[
 \operatorname{Sub}_n(T,\rho)
 =\{S\in\U_n\mid S\subseteq\sem T_\rho\}.
\]
A value binder $x\mathbin{\mathrm{in}}T$ ranges over $\sem T_\rho$.
Bounds may depend on earlier variables. Their context extension is the
comprehension set
\[
 \Gamma.D=\{(\rho,d)\mid\rho\in\Env(\Gamma),\ d\in D_\rho\}.
\]
Equations \eqref{eq:exists} and \eqref{eq:forall} then quantify over the fiber
$D_\rho$. This is the same construction for constant and dependent ranges.

For fixed ambient sorts, bounded quantification can be expressed by guards:
\[
 \exists a\in D\,P=\exists a\,(D(a)\land P),\qquad
 \forall a\in D\,P=\forall a\,(D(a)\Rightarrow P).
\]
The implication is a semantic operation in the Boolean predicate model; a
surface implementation can retain the guarded clause directly.

\section{Substitution and duality}
For a pullback square of context maps, reindexing commutes with both
quantifiers:
\[
 h^*(\exists_\pi P)=\exists_{\pi'}(h'^*P),\qquad
 h^*(\forall_\pi P)=\forall_{\pi'}(h'^*P).
\]
To verify either equation, fix the base assignment and compare the corresponding
fibers of the pullback. Their witnesses are in bijection. These Beck--Chevalley
laws govern capture-avoiding copying of quantified descriptions.

In each Boolean fiber,
\[
 \forall_\pi P=\neg\exists_\pi(\neg P).
\]
The semantic duality is available even in a surface fragment that exposes only
positive conjunctions, disjunctions, and universal binders.

\section{Empty ranges and uniform subjects}
The empty-index conventions are essential:
\[
 \bigcap\varnothing=V,\qquad\bigcup\varnothing=\varnothing.
\]
Hence $\forall a\in\varnothing\,P=\top$ and
$\exists a\in\varnothing\,P=\bot$. If $\U_n$ contains bottom, then
$\forall A\,A=\bot$. If lists are finite and bottom is admissible,
$\forall A\,\List(A)$ consists exactly of the empty list.

\example{One subject at every instance}{\Uone{}}
\par\noindent\begin{minipage}{\linewidth}
\begin{lstlisting}
empty(A): [...A]                 // []
impossible(A): A                 // _|_
impossibleRecord(n in 0 | 1): {
    value: n                     // one field cannot be both 0 and 1
}
\end{lstlisting}
\end{minipage}\par
These are intersections of predicates on one list, value, or record.
A parametric alias abbreviates a description by substitution. A
descriptor-producing function instead returns a description as a value.

\chapter{Unification, floating, and monotonicity}
\label{ch:refinement}
\section{Pointwise unification}
For a common range $D$,
\begin{equation}
 (\forall a\in D\,P(a))\land(\forall a\in D\,Q(a))
 =\forall a\in D\,(P(a)\land Q(a)).\label{eq:pointwise}
\end{equation}
This follows by regrouping intersections. Binder names are alpha-equivalent.
\example{Refining a polymorphic record across declarations}{\Uone{}}
\par\noindent\begin{minipage}{\linewidth}
\begin{lstlisting}
r(A): {wrap: func(A) -> {value: A}}
r(B): {wrap: func(B) -> {tag: "wrapped"}}

// Unified description, after alpha-renaming:
r(T): {wrap: func(T) -> {value: T, tag: "wrapped"}}
\end{lstlisting}
\end{minipage}\par
The two output obligations constrain one operation at each $T$.

Equal spelling does not identify independent existential witnesses:
\[
 (\exists a\,P(a))\land(\exists b\,Q(b))
 =\exists a,b\,(P(a)\land Q(b)).
\]
Shared existential witnesses are introduced by a shared scope and are conjoined
before their binder is discharged.

\section{Different bounds retain different obligations}
For a common body $P$,
\[
 (\forall a\in D\,P(a))\land(\forall a\in E\,P(a))
 =\forall a\in D\cup E\,P(a).
\]
Here $D\cup E$ is a union of \emph{assignments}. For type bounds,
\[
 \operatorname{Sub}_n(B)\cup\operatorname{Sub}_n(C)
 \subseteq\operatorname{Sub}_n(B\cup C)
\]
can be strict: mixed subtypes of $B\cup C$ need lie below neither arm.
\example{Two bounds do not become one upper bound}{\Uone{}}
\par\noindent\begin{minipage}{\linewidth}
\begin{lstlisting}
choose(A: int):    func(A, A) -> A
choose(A: string): func(A, A) -> A

// The stronger declaration also covers mixed unions:
stronger(A: int | string): func(A, A) -> A
\end{lstlisting}
\end{minipage}\par
The first pair retains both guards. With different bodies, its meaning is
$\forall A\,((A\subseteq\Int\Rightarrow P(A))\land
(A\subseteq\Str\Rightarrow Q(A)))$.

\section{Legal scope extrusion}
For $a\notin\FV(Q)$, existential Frobenius reciprocity gives
\[
 (\exists a\,P)\land Q=\exists a\,(P\land Q).
\]
For a nonempty fixed range, the universal counterpart holds:
\[
 (\forall a\,P)\land Q=\forall a\,(P\land Q).
\]
For a possibly empty dependent range, the universally quantified expression
must retain $Q$ outside the guard. Implementations store the guard and the
independent conjunct separately until inhabitation is known.

\example{An ordinary field floats as a constant obligation}{\Uone{}}
\par\noindent\begin{minipage}{\linewidth}
\begin{lstlisting}
r(A): {id: func(A) -> A}
r:    {label: "utilities"}

// Type universes are nonempty:
r(A): {id: func(A) -> A, label: "utilities"}
\end{lstlisting}
\end{minipage}\par
Independent same-kind quantifiers commute. An existential and a universal
commute only when a separately justified independence law applies.
In general,
\[
 \exists y\,\forall x\,(y=x)\quad\ne\quad
 \forall x\,\exists y\,(y=x).
\]
The dependency graph records which universal choices a witness may depend on.
File order is irrelevant; dependency order is semantic.

\section{Disjunction and projection}
Existential projection preserves unions; universal projection preserves
intersections. Their opposite distributivities fail.
\example{Binder scope across a union}{\Uone{}+\Dep{}}
\par\noindent\begin{minipage}{\linewidth}
\begin{lstlisting}
a: forall (n in 0 | 1) (n | (1 - n))  // 0 | 1
b: (forall (n in 0 | 1) n) |
   (forall (n in 0 | 1) (1 - n))      // _|_
\end{lstlisting}
\end{minipage}\par
Projection of an entire record must also retain joint witnesses. Take
$R_0=\{(s,0)\}$ and $R_1=\{(s,1)\}$. Projection $p$ onto the first coordinate
satisfies
\[
 p(R_0\cap R_1)=\varnothing,
 \qquad p(R_0)\cap p(R_1)=\{s\}.
\]
Pushing record selection through a universal intersection is therefore only an
inclusion in general. Pointwise unification operates on the whole subject,
including its correlated fields.

\section{Refinement theorem}
\begin{theorem}[Monotone quantified refinement]\label{thm:monotone}
Let $P'\subseteq P$ be predicates on one elaborated context, with the same
quantifier dependency graph. Every legal composite $\mathcal Q$ of existential
and universal projections satisfies $\mathcal Q(P')\subseteq\mathcal Q(P)$.
In particular, $\mathcal Q(P\land R)\subseteq\mathcal Q(P)$.
\end{theorem}
\begin{proof}
Both quantifiers are monotone, either by \cref{prop:adjoints} or directly from
their membership clauses. Composition preserves monotonicity.
\end{proof}
A contradictory closed predicate remains contradictory. A projected atomic
result constrained to $\{v\}$ can remain $\{v\}$ or become empty. Its successful
completion cannot change to a different atom.

\section{Refinable bounds are correlated variables}
Let an existential semantic type $X$ have upper bound $\Num$, and let
\[
 P(X,f)\equiv X\subseteq\Num\ \land\ \forall A\subseteq X\,T(A,f).
\]
Adding $X\subseteq\Int$ gives $P'=P\land(X\subseteq\Int)\subseteq P$.
The result remains monotone after existentially hiding $X$.

The universal obligation for one fixed $X$ is contravariant in its bound; the
\emph{joint predicate on possible $X$ and $f$} is refined by conjunction.
A solver retains $X$ and its correlations. It does not substitute its current
upper approximation into every occurrence.

\example{A particular client and a universally capable producer}{\Uone{}}
\par\noindent\begin{minipage}{\linewidth}
\begin{lstlisting}
job: {
    x: number
    f: func(x) -> string
    result: f(x)
}
job: {x: int & >=0}
job: {x: 7, f: func(n: int) -> string: "\(n)"}
\end{lstlisting}
\end{minipage}\par
In a complete instance, \code{x} denotes one integer and the function obligation
uses its singleton type. The joint relation retains that dependency. A separate
\code{func(number) -> string} requires all numeric inputs.
\todoob The elaborator distinguishes an ordinary witness reference, whose
type-position use is a singleton, from a type-sorted variable, whose use is its
decoded semantic type. Copying a CUE description preserves that classification
and the copied witness graph.

\chapter{Function types from universal obligations}
\label{ch:functions}
\section{Operational inhabitants and calls}
A closure is an operational value with code, a captured environment, and a fixed
argument-binding protocol. Let $F\subseteq V$ be the set of such values. A
\emph{call packet} records positional arguments, labeled arguments, and explicit
omission. A deterministic partial evaluator has observations
\[
 \Run(f,p)\in\{\Ok(S),\Reject,\Raise(e)\},
\]
where $S\subseteq V$ is a set of successful result completions. A concrete result
$v$ gives $S=\{v\}$. Evaluation can also remain unresolved; unresolved evaluation
is a state of computation, rather than a proof that its completion set is empty.
The terminating reference fragment uses structurally bounded evaluation.

Define the successful-call predicate
\[
 G_B(f,p)\equiv
 \exists S.\;\Run(f,p)=\Ok(S)\land\varnothing\ne S\subseteq B.
\]
A function contract is then
\begin{equation}
 A\Rightarrow B
 =\{f\in F\mid\forall p\in A.\;G_B(f,p)\}.\label{eq:arrow}
\end{equation}
Here $A$ is a set of complete call packets, derived from the parameter surface.
Single positional arguments can be identified with their one-element packets.
The surface \code{func(A) -> B} denotes this set.

A body supplies one operational inhabitant. Contract membership asserts its
coverage and result property. The predicate does not perform program synthesis.
For an empty domain, $\varnothing\Rightarrow B=F$. For nonempty $A$,
$A\Rightarrow\varnothing=\varnothing$ in the successful-call interpretation.

\section{Variance and intersections}
\begin{proposition}[Arrow laws]\label{prop:arrow}
For all domain and result sets:
\begin{align*}
 A\subseteq C,\ D\subseteq B
 &\Longrightarrow(C\Rightarrow D)\subseteq(A\Rightarrow B),\\
 (A\cup C)\Rightarrow B
 &=(A\Rightarrow B)\cap(C\Rightarrow B),\\
 A\Rightarrow(B\cap D)
 &=(A\Rightarrow B)\cap(A\Rightarrow D).
\end{align*}
\end{proposition}
\begin{proof}
The first two statements restrict or split the universal input obligation.
For the third, determinism fixes the same result description $S$ at each packet;
requiring $S\subseteq B$ and $S\subseteq D$ is equivalent to requiring
$S\subseteq B\cap D$, with the same success and nonemptiness conditions.
\end{proof}
Thus $(\Int\Rightarrow\Str)\cap(\Num\Rightarrow\Str)
=\Num\Rightarrow\Str$. More generally, an intersection of arrows preserves
input/output correlations that a single arrow may lose.

\example{A correlation-preserving overloaded contract}{\Uone{}}
\par\noindent\begin{minipage}{\linewidth}
\begin{lstlisting}
convert: (func(int) -> string) & (func(string) -> int)
\end{lstlisting}
\end{minipage}\par
One implementation must meet both obligations. A body may perform explicit
case analysis. The type intersection itself performs no branch selection.
On overlapping domains, all result obligations apply.

\section{Implementations and application refinement}
For a complete closure $f$, unification with $T$ denotes $\{f\}\cap\sem T$.
A successful conformance check preserves the closure. An incompatibility makes
that intersection empty. A symbolic closure carries unresolved environment and
conformance predicates until sufficient information is available.

For a set of closures $X$ and input packets $Q$, define
\[
 \App(X,Q)=\bigcup\{S\mid f\in X,\ p\in Q,\ \Run(f,p)=\Ok(S)\}.
\]
Joint environments are retained when $X$ and $Q$ share variables.
Then
\begin{equation}
 \App(X\cap Y,Q)\subseteq\App(X,Q)\cap\App(Y,Q).\label{eq:appmeet}
\end{equation}
Membership on the left has one closure witness that works on both sides.
Equality can fail: distinct closures can agree at a particular input while
having empty singleton intersection.

\example{Capabilities and one-call constraints}{\Uone{}}
\par\noindent\begin{minipage}{\linewidth}
\begin{lstlisting}
wide:   func(number) -> string
wide:   func(x: int) -> string: "ok"  // incompatible coverage

narrow: func(int) -> string
narrow: func(x: number) -> string: "ok" // conforming implementation
\end{lstlisting}
\end{minipage}\par
The second contract states a minimum capability. A particular call domain can
be constrained with an existential argument cell as in \cref{ch:refinement}.
This is the quantifier distinction between clients and producers.

\section{Erased polymorphic functions}
A polymorphic identity has the set interpretation
\[
 I=\bigcap_{A\in\U_n}(A\Rightarrow A).
\]
If all relevant singletons are admissible, then $f\in I$ implies
$\Run(f,x)=\Ok(\{x\})$ for each $x$: instantiate at $\{x\}$.
If the admissible family is coarser, the statement only forces preservation of
its available predicates.

\example{A scheme and an ordinary arrow unify normally}{\Uone{}}
\par\noindent\begin{minipage}{\linewidth}
\begin{lstlisting}
id(A): func(A) -> A
id:    func(int) -> int       // redundant instance

bad(A): func(A) -> A
bad:    func(int) -> bool     // _|_: int and bool are disjoint
\end{lstlisting}
\end{minipage}\par
Universals distribute through these conjunctions by \cref{eq:pointwise}. An
ordinary instantiation is a consequence of the universal constraint, rather
than its replacement.

\section{Checked calls and effects}
For a set $\epsilon$ of admitted error tags, define
\[
 A\Rightarrow_\epsilon B
 =\{f\in F\mid\forall p\in A.\;
       G_B(f,p)\ \lor\ \exists e\in\epsilon.\;\Run(f,p)=\Raise(e)\}.
\]
The pure arrow has $\epsilon=\varnothing$. Domain rejection is a separate
observation from a permitted failure within the domain. A bridge supporting
logical integers with representability failures has type
$\Int\Rightarrow_{\mathsf{bridge}}\Int$. It does not thereby support arbitrary
fractional inputs. Divergence, when admitted by an execution profile, is a
separately named outcome.


\chapter{Surface binding and conservative extension}
\label{ch:surface}
\section{Quantifier expressions}
The core surface expressions are
\par\noindent\begin{minipage}{\linewidth}
\begin{lstlisting}
forall (A, B) T
exists A { ... }
forall (A: number) func(A, A) -> [A, A]
forall (n in int & >=0) T
\end{lstlisting}
\end{minipage}\par
Parentheses group multiple binders or a constrained binder; a single simple
binder may omit them. The body is the following expression. A record body is
already delimited by braces. Quantifiers have low expression precedence, and
parentheses delimit narrower scopes inside unions and arrows.

Unadorned $A$ is type-sorted, at an inferred universe. \code{A: T} is an upper
bound on that semantic type. \code{x in T} binds an ordinary value. An explicit
\code{A in Type(n): T} specifies both the universe and the upper bound.
All binders are lexical and alpha-renamable. Bounds may reference earlier
binders and outer variables.

\section{Quantified field declarations}
The declaration
\par\noindent\begin{minipage}{\linewidth}
\begin{lstlisting}
map(A, B): func(func(A) -> B, [...A]) -> [...B]
\end{lstlisting}
\end{minipage}\par
elaborates exactly to
\par\noindent\begin{minipage}{\linewidth}
\begin{lstlisting}
map: forall (A, B) func(func(A) -> B, [...A]) -> [...B]
\end{lstlisting}
\end{minipage}\par
The field value is fixed outside the quantifier. The outer record's field
selection is retained as the subject of the quantified predicate.

For a quantified record:
\par\noindent\begin{minipage}{\linewidth}
\begin{lstlisting}
r(A): {
    id:   func(A) -> A
    wrap: func(A) -> {value: A}
}
\end{lstlisting}
\end{minipage}\par
the order is $\exists r\,\forall A\,P(r,A)$ when the surrounding record is
closed existentially. Both projections refer to the same record for all $A$.
The field-sugar form is concise for rank-1 declarations; RHS syntax also permits
nested and higher-rank positions.

\section{Quantifier blocks}
A lexical binder prefix can be written in a record block:
\par\noindent\begin{minipage}{\linewidth}
\begin{lstlisting}
x: {
    exists T: U
    forall V: W
    // The whole following declaration body is in this prefix.
    value: T
    transform: func(V) -> V
}
\end{lstlisting}
\end{minipage}\par
This elaborates to
\par\noindent\begin{minipage}{\linewidth}
\begin{lstlisting}
x: exists (T: U) forall (V: W) {
    value: T
    transform: func(V) -> V
}
\end{lstlisting}
\end{minipage}\par
The ordered prefix records dependency. The declarations after the prefix are
conjoined without source-order dependence. Contributions from other files
unify with the entire quantified subject; they do not capture these private
binder names merely by spelling them alike. A named interface may explicitly
export a sharing equation when such sharing is intended.

A syntactically asymmetric alternative requires external parentheses for
universals and block prefixes for existentials:
\par\noindent\begin{minipage}{\linewidth}
\begin{lstlisting}
x(V: W): {
    exists T: U
    value: T
    transform: func(V) -> V
}
\end{lstlisting}
\end{minipage}\par
Its prefix is $\forall V\,\exists T$, rather than $\exists T\,\forall V$.
The two spellings express different dependence. Both quantifiers still use the
dual semantic operations of \cref{ch:quantifiers}.

\section{Function-local generics and parametric aliases}
A compact function-only spelling is possible:
\par\noindent\begin{minipage}{\linewidth}
\begin{lstlisting}
map: func<A, B>(func(A) -> B, [...A]) -> [...B]
\end{lstlisting}
\end{minipage}\par
It expands to an outer universal on the function subject.

A parametric alias binds parameters in a description:
\par\noindent\begin{minipage}{\linewidth}
\begin{lstlisting}
Box(A) = {value: A}
Pair(A, B) = [A, B]
\end{lstlisting}
\end{minipage}\par
The \code{=} form defines an alias; the \code{:} form constrains a field.
Alias application substitutes its arguments capture-avoidingly into the body,
so \code{Box(int)} expands to \code{{value: int}}. Alias parameters use the
same sort and bound syntax as quantified binders; arguments must satisfy those
bounds. A nonparametric alias uses the existing form \code{let S = T}.

In contrast, \code{box(A): {value: A}} requires one value lying in every $A$
and is bottom when bottom is admissible. Reified descriptor functions retain
their descriptor arguments as correlated variables.

\code{foo[T]} selects the instance of a universally constrained subject
justified by substituting \code{T} for its binder. For a universal description,
selection refers to the same inhabitant at that instance; it retains the subject
and its universal constraints. Thus \code{id[int](3)} uses the integer instance
of \code{id}. Selection is erased. Implicit instantiation uses fresh flexible
variables at the call site and retains the universal declaration.

\section{A conservative embedding}
\begin{definition}[Legacy-data observation]
A legacy-data observation is evaluation and export using the pre-extension
operators and data formats, applied to a program in the legacy syntax. Its
completion domain is $D$, embedded into the extended $V$ by $i:D\hookrightarrow V$.
\end{definition}

\Needspace{10\baselineskip}
\begin{theorem}[Relative conservativity]\label{thm:conservative}
Suppose the extension retains the legacy elaborator, base evaluation rules,
default-selection rules, and data constructors. For every legacy description
$T$ and legacy assignment $\rho$,
\[
 i^{-1}(\sem T^{+}_{i\rho})=\sem T^{0}_{\rho}.
\]
Closed legacy programs have the same legacy-data observations under the extended
evaluator.
\end{theorem}
\begin{proof}
For data predicates, the claim is the defining restriction of each extended
primitive to $D$. Inverse image preserves intersections and unions. Record and
list constructors preserve the embedding and the field-sharing elaboration.
Existing local projection and copying operations are unchanged. The operational
statement follows by induction on the unchanged evaluation derivation, with
identical default selection and export.
\end{proof}
The theorem is observational and data-relative. Extended top additionally
admits new kinds of inhabitants in new-language contexts. This distinguishes
conservativity from equality of global carriers.

\decision{Gate new binder and arrow forms by language version or experiment
attributes. Treat \code{forall}, \code{exists}, and \code{func} as contextual
keywords in the selected syntax. Existing data files retain their parsing,
closing, copying, and evaluation behavior.}

Existing struct-as-function code remains a relational program:
\par\noindent\begin{minipage}{\linewidth}
\begin{lstlisting}
#Add: {a: int, b: int, out: a + b}
sum: (#Add & {a: 1, b: 2}).out
\end{lstlisting}
\end{minipage}\par
A function declaration adds a reusable universal obligation and a concise call
form. Refactoring to that declaration is an explicit program change, checked
against the resulting contract.

\chapter{Three implementation profiles}
\label{ch:profiles}
\section{Rank, predicativity, and interface boundaries}
Type rank measures the placement of polymorphic binders relative to function
arguments. A callback required to work at all types is higher-rank; an ordinary
higher-order function such as \code{map} can have a rank-1 type. Predicativity is
a separate restriction on the universes admitted as instantiations.

For the restricted profiles, use a \emph{boundary-prenex} discipline. Within one
executable signature clause, type quantifiers occur in an outer prefix over a
quantifier-free body, after justified normalization. Module interfaces are
explicit additional boundaries. Finite intersections retain several prefixed
clauses when their bounds or dependencies differ. An alias expands for the rank
check; naming a higher-rank type does not change its rank.

A positive universal may sometimes float through a result arrow when its
variable is absent from the input and the established deterministic arrow law
justifies the move. Quantifiers over argument types retain their position.
Existentials whose witnesses depend on arguments likewise retain that dependence.

\section{Full symmetric, higher-rank CUE: \Shigher{}}
\Shigher{} admits both quantifiers at arbitrary well-sorted type positions, with the
predicative universe discipline. It includes first-class existential packages,
polymorphic arguments, and nested quantified records.
\par\noindent\begin{minipage}{\linewidth}
\begin{lstlisting}
use: func(forall A func(A) -> A) -> [int, string]

#Showable: exists A {
    value: A
    show:  func(A) -> string
}
show: func(#Showable) -> string
\end{lstlisting}
\end{minipage}\par
The checker uses explicit introduction and elimination rules, scope levels, and
annotations. Higher-rank bidirectional checking provides a useful proof-theoretic
organization \cite{dunfield}, while Boolean connectives and dependent predicates
require their own sound constraint procedures.

\section{Symmetric, rank-1 CUE: \Sone{}}
\Sone{} admits existential and universal prefixes at declaration and module-interface
boundaries. An executable callback parameter is monomorphic at each call.
Abstract modules are opened at their module boundary, exposing a rigid local
representation name and monomorphic operations to their clients. Passing an
inline quantified package to a function, or receiving a universally quantified
callback, uses \Shigher{}.

A representative interface is
\par\noindent\begin{minipage}{\linewidth}
\begin{lstlisting}
Counter: exists State {
    zero: State
    next: func(State) -> State
    read: func(State) -> int
}
\end{lstlisting}
\end{minipage}\par
Module-level client code opens \code{Counter}, then passes \code{next} and
\code{read} as ordinary monomorphic functions. The prefix can contain both kinds
of binder, with explicit dependency order. Symmetry concerns their availability
and semantic laws at the same boundaries; it does not authorize commuting them.

This profile supports ML-style abstract module reuse with a boundary-prenex
rank restriction.

\section{Universal-only, rank-1 CUE: \Uone{}}
\Uone{} adds only explicit outer universal type binders. Existing instance variables,
local field cells, and call-instantiation variables remain existentially
interpreted. The surface need not expose existential type packaging to express
identity, mapping, composition, filtering, folds supplied as primitives, and
many generic adapter interfaces.

The same semantics still defines existential projection internally. The
implementation simply exposes a smaller formation and elimination interface.
It therefore gains universal capability obligations and compositional
polymorphic subtyping without requiring an opaque-module boundary in its first
release.

\par\noindent\begin{minipage}{\linewidth}
\begin{lstlisting}
id(A): func(A) -> A
compose(A, B, C): func(func(B) -> C, func(A) -> B) ->
    func(A) -> C
map(A, B): func(func(A) -> B, [...A]) -> [...B]
\end{lstlisting}
\end{minipage}\par
These declarations require rank-1 universals. They already distinguish reusable
implementations from existentially constrained calls.

\section{Coverage of host-language interfaces}
Rank-1 signatures describe ordinary generic APIs. Higher-rank callbacks,
abstract packages, and some lifetime-polymorphic or dependent APIs require a
richer representation.
GHC's documented distinction between rank-1 schemes and polymorphic parameters
is a direct example \cite{ghcrank}. A host binding must either express its actual
quantifier structure or export a deliberately weaker interface with that loss
recorded.

\begin{tabularx}{\textwidth}{@{}Yccc@{}}
\toprule
Facility & \Uone{} & \Sone{} & \Shigher{}\\\midrule
Outer polymorphic arrows and records & yes & yes & yes\\
Finite intersections of scheme clauses & yes & yes & yes\\
Existential instance variables & yes & yes & yes\\
Explicit abstract module type witnesses & --- & yes & yes\\
Polymorphic callback parameters & --- & --- & yes\\
First-class nested quantified packages & --- & --- & yes\\
Value dependence as a separate extension & \Dep{} & \Dep{} & \Dep{}\\
\bottomrule
\end{tabularx}

\section{Implementation choice}
\Uone{} supports rank-1 generic interfaces. \Sone{} adds explicit modular abstraction
at stable boundaries. \Shigher{} admits quantified descriptions at arbitrary
well-sorted positions. The inclusion of source fragments is conservative on
programs already admitted by the smaller profile.
A rank restriction does not by itself decide arbitrary semantic subtyping,
quantified arithmetic, or program equivalence; each checker specifies its
supported predicate fragment.

\chapter{Programs and calling conventions}
\label{ch:programs}
\section{Basic polymorphic constructions}
\example{Identity at unrelated instances}{\Uone{}}
\par\noindent\begin{minipage}{\linewidth}
\begin{lstlisting}
id(A): func(x: A) -> A: x
integer: id(7)                 // 7
text:    id("seven")           // "seven"
record:  id({port: 443})       // {port: 443}
chosen:  id[int & >=0](7)      // a checked explicit instance
\end{lstlisting}
\end{minipage}\par
One closure meets every instance. Explicit instantiation does not construct a
wrapper or remove other valid instances.

\example{Pairing preserves an inferred common refinement}{\Uone{}}
\par\noindent\begin{minipage}{\linewidth}
\begin{lstlisting}
pair(A: number): func(x: A, y: A) -> [A, A]: [x, y]
p: pair[int & >=1](2, 5)       // [2, 5]
q: pair(1, 2.5)                // [1, 2.5]
\end{lstlisting}
\end{minipage}\par
The mixed call has the admissible instance $A=\{1,2.5\}$. A shared type variable
preserves a common refinement; it does not enforce equal primitive tags.
A constant result $[0,0]$ would satisfy the unrefined numeric result type but
would fail the instance $A=\{2,5\}$. Matching primitive categories can instead
be specified by the input relation $[\Int,\Int]\cup[\mathsf{Float},\mathsf{Float}]$.

\example{Selecting a value preserves its subtype}{\Uone{}}
\par\noindent\begin{minipage}{\linewidth}
\begin{lstlisting}
min(A: number): func(x: A, y: A) -> A: {
    if x <= y {out: x}
    if x > y  {out: y}
}.out

m: min[int & >=10](12, 20)     // 12
\end{lstlisting}
\end{minipage}\par
Either branch returns an input inhabitant. By contrast, the proposed contract
\code{forall (A: number) func(A) -> A} does not justify adding one: $A=\{0\}$
is a counterexample. Numeric bounds permit numeric operations; result
refinements require their own proof.

\example{Two type parameters preserve asymmetric information}{\Uone{}}
\par\noindent\begin{minipage}{\linewidth}
\begin{lstlisting}
swap(A, B): func(p: [A, B]) -> [B, A]: [p[1], p[0]]
s: swap([7, "seven"])         // ["seven", 7]
\end{lstlisting}
\end{minipage}\par
The two coordinates retain their distinct type variables, rather than being
merged into one union.

\Needspace{180pt}
\section{Generic collection programs}
\example{Map with a single implementation}{\Uone{}}
\par\noindent\begin{minipage}{\linewidth}
\begin{lstlisting}
map(A, B): func(g: func(A) -> B, xs: [...A]) -> [...B]: [
    for x in xs {g(x)}
]
inc:  func(x: int) -> int: x + 1
text: func(x: int) -> string: "\(x)"

numbers: map(inc, [1, 2, 3])    // [2, 3, 4]
strings: map(text, [1, 2, 3])   // ["1", "2", "3"]
\end{lstlisting}
\end{minipage}\par
The body is checked with rigid $A,B$. Each call creates flexible instance
variables. Finite sets of admissible instantiations can contribute intersected
typings, as in the CDuce polymorphic application discipline \cite{cducetutorial}.

\example{Filtering and partitioning preserve elements}{\Uone{}}
\par\noindent\begin{minipage}{\linewidth}
\begin{lstlisting}
filter(A): func(p: func(A) -> bool, xs: [...A]) -> [...A]: [
    for x in xs if p(x) {x}
]
partition(A): func(p: func(A) -> bool, xs: [...A]) -> {
    yes: [...A]
    no:  [...A]
}: {
    yes: [for x in xs if p(x) {x}]
    no:  [for x in xs if !p(x) {x}]
}
positive: func(x: int) -> bool: x > 0
kept: filter(positive, [-2, 0, 4]) // [4]
\end{lstlisting}
\end{minipage}\par
The callback is pure. An effectful callback requires an explicit effect policy,
including the meaning of repeated evaluation.

\example{Flattening a generic transformation}{\Uone{}}
\par\noindent\begin{minipage}{\linewidth}
\begin{lstlisting}
flatMap(A, B): func(g: func(A) -> [...B], xs: [...A]) -> [...B]: [
    for x in xs for y in g(x) {y}
]
duplicate(A): func(x: A) -> [A, A]: [x, x]
twice: flatMap(duplicate, [1, 2]) // [1, 1, 2, 2]
\end{lstlisting}
\end{minipage}\par
\example{Mapping named record coordinates}{\Uone{}}
\par\noindent\begin{minipage}{\linewidth}
\begin{lstlisting}
mapPoint(A, B): func(g: func(A) -> B, p: {x: A, y: A}) -> {
    x: B
    y: B
}: {
    x: g(p.x)
    y: g(p.y)
}
p: mapPoint(text, {x: 2, y: 3}) // {x: "2", y: "3"}
\end{lstlisting}
\end{minipage}\par
Open/closed record constraints retain their existing meaning. A separate row
parameter can preserve additional fields when a row-polymorphic fragment is
selected.

\section{Higher-order composition and return values}
\example{Composition is rank-1}{\Uone{}}
\par\noindent\begin{minipage}{\linewidth}
\begin{lstlisting}
compose(A, B, C): func(g: func(B) -> C, f: func(A) -> B) ->
    func(A) -> C:
    func(x: A) -> C: g(f(x))

pipeline: compose(text, inc)
answer: pipeline(4)           // "5"
\end{lstlisting}
\end{minipage}\par
The returned closure captures $f$ and $g$. Its generic variables are rigid
assumptions while checking the body; their run-time representations are erased.

\example{A value-dependent closure}{\Uone{}}
\par\noindent\begin{minipage}{\linewidth}
\begin{lstlisting}
withPrefix: func(prefix: string) -> func(string) -> string:
    func(s: string) -> string: prefix + s

warn: withPrefix("warning: ")
message: warn("retry")        // "warning: retry"
\end{lstlisting}
\end{minipage}\par
Two closures from the same body with different concrete prefixes have different
captured environments.

\example{A polymorphic callback}{\Shigher{}}
\par\noindent\begin{minipage}{\linewidth}
\begin{lstlisting}
useBoth: func(p: forall A func(A) -> A) -> [int, string]: [
    p(7),
    p("seven"),
]
ok: useBoth(id)              // [7, "seven"]
onlyInt: func(x: int) -> int: x
bad: useBoth(onlyInt)         // incompatible universal obligation
\end{lstlisting}
\end{minipage}\par
The callback itself is polymorphic. A new instantiation of the surrounding
function for each use of $p$ would not prove this contract.

\example{A nested universal result}{\Shigher{} syntax; prenex normalization may apply}
\par\noindent\begin{minipage}{\linewidth}
\begin{lstlisting}
constant(A): func(x: A) -> (forall B func(B) -> A):
    func(_) -> A: x

seven: constant(7)
a: seven("ignored")          // 7
b: seven(true)               // 7
\end{lstlisting}
\end{minipage}\par
The resulting closure depends on $x$ and works uniformly for every subsequent
argument type. Moving $B$ outward is permitted only through the proven positive
arrow law, with $B$ absent from the domain.

\section{Quantified records and independent refinement}
\example{A reusable polymorphic utility record}{\Uone{}}
\par\noindent\begin{minipage}{\linewidth}
\begin{lstlisting}
utilities(A): {
    id:   func(x: A) -> A: x
    wrap: func(x: A) -> {value: A}: {value: x}
}
utilities: {name: "core"}

wrapped: utilities.wrap("payload") // {value: "payload"}
\end{lstlisting}
\end{minipage}\par
The \code{name} field is a constant conjunct on the same record.

\example{A cross-file behavioral strengthening}{\Uone{}}
\par\noindent\begin{minipage}{\linewidth}
\begin{lstlisting}
// interface.cue
wrap(A): func(A) -> {value: A}

// metadata.cue
wrap(A): func(A) -> {tag: "wrapped"}

// implementation.cue
wrap(A): func(x: A) -> {value: A, tag: "wrapped"}: {
    value: x
    tag: "wrapped"
}
\end{lstlisting}
\end{minipage}\par
The body is checked against both pointwise result constraints. No declaration
replaces or takes precedence over another.

\example{An intersection of domain-specific guarantees}{\Uone{}}
\par\noindent\begin{minipage}{\linewidth}
\begin{lstlisting}
classify: func(x: int) -> ("negative" | "nonnegative"): {
    if x < 0  {out: "negative"}
    if x >= 0 {out: "nonnegative"}
}.out
classify: (func(int & <0) -> "negative") &
          (func(int & >=0) -> "nonnegative")

c: classify(-3)              // "negative"
\end{lstlisting}
\end{minipage}\par
An ordinary union of function values describes alternatives, whereas this
intersection describes one value satisfying both implications.

\example{Unresolved alternatives and defaults}{\Uone{}}
\par\noindent\begin{minipage}{\linewidth}
\begin{lstlisting}
f: (func(x: int) -> int: x) |
   (func(x: int) -> int: x + 1)
a: f(3)                     // incomplete operational choice

preferred: *(func(x: int) -> int: x) |
            (func(x: int) -> int: x + 1)
b: preferred(3)             // 3, under legacy default selection
\end{lstlisting}
\end{minipage}\par
Constraint propagation can retain the set of alternative completions. Execution
uses the ordinary concreteness/default discipline to select a callable value.

\section{Argument forms and call packets}
The five parameter forms in proposal 4484 all elaborate to packet predicates
\cite{cueproposal}. Named parameters admit a positional packet or the matching
label; required and optional name-only parameters use labels; anonymous and
aliased positional parameters use positions only. Packet binding precedes value
constraint solving.

\example{Positional, named, and mixed calls}{\Uone{}}
\par\noindent\begin{minipage}{\linewidth}
\begin{lstlisting}
sum: func(a: int, b: int) -> int: a + b
p: sum(2, 3)                // 5
q: sum(a: 2, b: 3)          // 5
r: sum(2, b: 3)             // 5
\end{lstlisting}
\end{minipage}\par

\example{Name-only required and optional parameters}{\Uone{}}
\par\noindent\begin{minipage}{\linewidth}
\begin{lstlisting}
key: func(a!: int, note?: string) -> int: a
x: key(a: 4)                // 4; note is absent
z: key(a: 4, note: "audit") // 4
bad: key(4)                 // invalid packet: a is name-only
\end{lstlisting}
\end{minipage}\par
A body that reads an absent optional parameter must prove presence or produce a
checked absence failure. Optionality is an absence alternative, not an unknown
present value.

\example{Positional-only aliases and anonymous parameters}{\Uone{}}
\par\noindent\begin{minipage}{\linewidth}
\begin{lstlisting}
scale: func(_~value: int, _: int) -> int: value * 2
x: scale(5, 99)             // 10
bad: scale(value: 5, 99)     // invalid label/order
\end{lstlisting}
\end{minipage}\par
The positional alias names a local body variable without adding a public
label. Renaming it preserves the packet protocol.

\example{Hidden labels and attributes}{\Uone{}}
\par\noindent\begin{minipage}{\linewidth}
\begin{lstlisting}
package internal
measure: func(_value: int @go(Int), unit!: string) -> int: _value
inside: measure(_value: 7, unit: "ms")
// Another package can use the positional slot, not the hidden label.
\end{lstlisting}
\end{minipage}\par
Hidden labels include their package identity. Attributes are retained as
metadata; only a declared extension assigns them operational significance.

\example{Errors distinguish presence from values}{\Uone{}}
\par\noindent\begin{minipage}{\linewidth}
\begin{lstlisting}
a: sum(1, a: 1)            // duplicate binding, even with equal values
b: sum(1, 2, 3)            // surplus positional argument
c: sum(a: 1, b: 2, c: 3)    // unknown label
d: sum(1)                  // required argument missing
e: sum(_, 2)               // supplied but incomplete, not missing
\end{lstlisting}
\end{minipage}\par
A duplicate packet is invalid before field unification. Conjoining equal values
does not erase a duplicate-binding error.

\section{Defaults}
\example{Omission defaults and caller preferences}{\Uone{}}
\par\noindent\begin{minipage}{\linewidth}
\begin{lstlisting}
add: func(a: int, b: int = 10) -> int: a + b
x: int | *5
omitted:  add(1)            // 11
supplied: add(1, x)         // 6
unknown:  add(1, _)         // incomplete; b was supplied
\end{lstlisting}
\end{minipage}\par
An omission default belongs to the closure protocol. A supplied argument bypasses
it and retains its own ordinary preference information. The default expression
is checked against its parameter description in the declaration environment.

\example{A defaulted name-only argument}{\Uone{}}
\par\noindent\begin{minipage}{\linewidth}
\begin{lstlisting}
render: func(text: string, suffix!: string = "!") -> string:
    text + suffix
x: render("hello")                 // "hello!"
y: render("hello", suffix: "?")   // "hello?"
\end{lstlisting}
\end{minipage}\par
Defaulted required-label parameters may be omitted, while optional parameters
retain absence semantics and have no omission default.

\example{A configurable default is an explicit adapter}{\Uone{}}
\par\noindent\begin{minipage}{\linewidth}
\begin{lstlisting}
raw: func(x: int) -> int: x
configured: func(x: int = 2) -> int: raw(x)
answer: configured()        // 2
\end{lstlisting}
\end{minipage}\par
The adapter supplies the protocol that accepts the empty packet. A contract
mentioning that protocol requires an implementation possessing it. It does not
mutate a previously supplied closure.

\section{Open interfaces and partial application}
An open signature describes a partial \emph{protocol interface}. Elaborate its
ellipsis to an existential row of further parameter declarations, retaining
that witness with the function. Its callable packet domain is determined only
when the required row is known. This is different from promising that every
completion of the row can be omitted. \Uone{} retains this row as an internal
existential witness, just as it retains call-instantiation variables; explicit
existential type packaging is a separate surface facility.

\example{An open interface completed by an implementation}{\Uone{}, implicit protocol row}
\par\noindent\begin{minipage}{\linewidth}
\begin{lstlisting}
T: func(a: int, ...) -> number
f: T & (func(a: int, b: int) -> int: a + b)
x: f(1, 2)                 // 3
y: f(1)                    // missing b
\end{lstlisting}
\end{minipage}\par
The row witness is completed by \code{b: int}. A closed one-argument capability
admits an implementation with additional omittable parameters because its
promised one-argument packets remain covered.

\example{Chained partial application}{\Uone{}}
\par\noindent\begin{minipage}{\linewidth}
\begin{lstlisting}
add3: func(a: int, b: int, c: int) -> int: a + b + c
f: add3(1, ...)
g: f(2, ...)
x: g(3)                    // 6
y: add3(a: 1, ...)(2, 3)    // 6
\end{lstlisting}
\end{minipage}\par
A partial closure stores bindings to original parameter slots. Supplied values
are checked immediately; the body runs on completion. An unbound default is
preserved and a bound argument suppresses that default.

\example{A residual contract on a partial closure}{\Uone{}}
\par\noindent\begin{minipage}{\linewidth}
\begin{lstlisting}
f: add3(1, ...)
f: func(int, int) -> int
x: f(2, 3)                 // 6
\end{lstlisting}
\end{minipage}\par
The residual protocol is a normal callable contract. Normalizing the original
slot bindings makes partial applications comparable independently of whether
arguments were supplied by label or position.

\section{Identity, captured values, and result combination}
Use a closure descriptor $(o,\gamma,\beta)$: declaration origin $o$, captured
free environment $\gamma$, and normalized partial bindings $\beta$.
The origin survives copying; irrelevant ambient fields are absent from
$\gamma$. Symbolic captures remain constraint variables.

\example{Self-unification after copying and embedding}{\Uone{}}
\par\noindent\begin{minipage}{\linewidth}
\begin{lstlisting}
x: {f: func(y: int) -> int: y}
a: x.f & x.f
b: x.f & {x}.f
c: x.f & (x & {g: 1}).f
r: [a(2), b(2), c(2)]       // [2, 2, 2]
\end{lstlisting}
\end{minipage}\par

\example{Equal source, unequal captures}{\Uone{}}
\par\noindent\begin{minipage}{\linewidth}
\begin{lstlisting}
d: [for v in [1, 2] {func(y: int) -> int: y + v}]
e: d[0] & d[1]             // _|_: captures differ
\end{lstlisting}
\end{minipage}\par
For one origin, captured-environment constraints unify structurally. Distinct
known origins denote distinct intensional values. This supplies idempotence
without deciding equality of mathematical functions.

\example{Combining result descriptions explicitly}{\Uone{}}
\par\noindent\begin{minipage}{\linewidth}
\begin{lstlisting}
left:  func(x: int) -> {a: int}: {a: x}
right: func(x: int) -> {b: int}: {b: x}
both: func(x: int) -> {a: int, b: int}: left(x) & right(x)
r: both(3)                  // {a: 3, b: 3}
\end{lstlisting}
\end{minipage}\par
The body is a new implementation whose result is the CUE unification of the two
result descriptions. Totality requires their intersection to be inhabited at
every admitted input. A version with arbitrary possibly conflicting callbacks
uses a checked conflict effect.

\section{Validators and iteration}
\example{A predicate used as an ordinary constraint}{\Uone{}}
\par\noindent\begin{minipage}{\linewidth}
\begin{lstlisting}
positive: func(x: int) -> bool: x > 0
#Positive: {value: int, _valid: true & positive(value)}
n: (#Positive & {value: 7}).value   // 7
bad: #Positive & {value: -1}       // _|_
\end{lstlisting}
\end{minipage}\par
A validator is a predicate on an existential subject. It can be expressed with
existing struct machinery once predicate functions are available.

\example{Portable foreign contract}{\Uone{} with checked effects}
\par\noindent\begin{minipage}{\linewidth}
\begin{lstlisting}
externF: extern func(int) -> int !bridge
// A native-range failure is a bridge outcome.
// A fractional input is outside the logical domain.
\end{lstlisting}
\end{minipage}\par
The bridge checks logical kind, representation, conversion, and returned
values as specified by the source proposal \cite{cueaddendum}. The effect
annotation records those failures; it introduces no additional conversion
stage. Tool-level effects execute in an explicit sequencing context, while
constraint unification operates on their descriptions.

\example{Nested calls remain ordinary evaluation}{\Uone{}}
\par\noindent\begin{minipage}{\linewidth}
\begin{lstlisting}
twice: func(n: int) -> int: n + n
value: twice(twice(twice(2))) // 16
\end{lstlisting}
\end{minipage}\par
The initial execution profile preserves the host cycle discipline. A structurally
recursive extension can admit calls justified by a decreasing finite measure.
This choice is independent of quantifier rank.

\example{A structurally recursive fold}{\Uone{} with structural recursion}
\par\noindent\begin{minipage}{\linewidth}
\begin{lstlisting}
fold(A, B): func(step: func(B, A) -> B, seed: B, xs: [...A]) -> B: {
    if len(xs) == 0 {out: seed}
    if len(xs) > 0 {out: fold(step, step(seed, xs[0]), xs[1:])}
}.out
plus: func(x: int, y: int) -> int: x + y
sum: fold(plus, 0, [1, 2, 3]) // 6
\end{lstlisting}
\end{minipage}\par
Each recursive call reduces list length. The termination proof is independent
of polymorphism; a nonrecursive profile may supply the same contract as a trusted
primitive.

\chapter{Dependent descriptions and functions}
\label{ch:dependent}
\section{Dependent contexts}
A dependent context extends $\Gamma$ by a value $x$ satisfying a description
$A$ in that context:
\[
 \Env(\Gamma,x:A)
 =\{(\rho,x)\mid\rho\in\Env(\Gamma),\ x\in\sem A_\rho\}.
\]
The existential and universal adjoints already apply to its projection.
Dependence therefore changes the fibers over which quantification ranges,
rather than introducing another interpretation of refinement.

A dependent function contract uses its actual argument in the result predicate:
\[
 \Pi_{x:A}B(x)
 =\{f\in F\mid\forall x\in A.\;G_{B(x)}(f,x)\}.
\]
A dependent pair retains the witness:
\[
 \Sigma_{x:A}B(x)=\{(x,y)\mid x\in A,\ y\in B(x)\}.
\]
Logical $\forall x\,B(x)$ constrains \emph{one subject} at every $x$; a dependent
function instead relates inputs to results. Logical $\exists x\,B(x)$ projects
away $x$; a dependent pair retains it. The distinction is determined by the
subject and the application or pairing relation.

\decision{Within a dependent function signature, parameters enter scope from
left to right. A result constraint can refer to all parameters. Default
expressions can depend only on parameters whose binding is already established.
The dependency graph determines evaluation order; it is independent of the
source order of unrelated record declarations.}

\example{A directly dependent result}{\Uone{}+\Dep{}}
\par\noindent\begin{minipage}{\linewidth}
\begin{lstlisting}
successor: func(x: int) -> (x + 1): x + 1

// Equivalent pointwise specification:
successor(n in int): func(n) -> (n + 1)
\end{lstlisting}
\end{minipage}\par
The second form quantifies over input singletons. The first uses a dependent
result predicate on the actual argument.

\example{A dependent parameter constraint}{\Uone{}+\Dep{}}
\par\noindent\begin{minipage}{\linewidth}
\begin{lstlisting}
interval: func(lo: int, hi: int & >=lo) -> {
    low: lo
    high: hi
}: {low: lo, high: hi}

ok: interval(3, 7)          // {low: 3, high: 7}
bad: interval(7, 3)         // domain contradiction
\end{lstlisting}
\end{minipage}\par
The admitted domain is a relation on two arguments. It is not the Cartesian
product of their marginal types.

\Needspace{180pt}
\section{Indexed finite sequences}
The following aliases use the list validators supplied by CUE
\cite{cuelist}:
\par\noindent\begin{minipage}{\linewidth}
\begin{lstlisting}
import "list"

let Nat = int & >=0
Vec(A, n in Nat) = [...A] &
    list.MinItems(n) & list.MaxItems(n)
Fin(n in Nat) = int & >=0 & <n
\end{lstlisting}
\end{minipage}\par
Their reference semantics is
\[
 \sem{\operatorname{Vec}(A,n)}
   =\{xs\in\List(A)\mid |xs|=n\},\qquad
 \sem{\operatorname{Fin}(n)}=\{0,\ldots,n-1\}.
\]
A symbolic index remains a constraint variable. The checker uses the length
predicate and its lemmas; concrete validator execution is one implementation
strategy for ground indices.

\example{Length-preserving map}{\Uone{}+\Dep{}}
\par\noindent\begin{minipage}{\linewidth}
\begin{lstlisting}
map(A, B, n in Nat): func(
    g: func(A) -> B,
    xs: Vec(A, n),
) -> Vec(B, n): [for x in xs {g(x)}]
\end{lstlisting}
\end{minipage}\par
The universal $n$ records a property of the input, rather than supplying a
runtime integer. Its value is available as \code{len(xs)} when evaluation
requires it.

\example{Safe indexing, including the zero-length case}{\Uone{}+\Dep{}}
\par\noindent\begin{minipage}{\linewidth}
\begin{lstlisting}
index(A, n in Nat): func(xs: Vec(A, n), i: Fin(n)) -> A: xs[i]

x: index(["a", "b"], 1)    // "b"
bad: index([], 0)           // 0 is outside Fin(0)
\end{lstlisting}
\end{minipage}\par
For $n=0$, the packet domain is empty. The generic implementation remains valid;
there is no admitted out-of-range call.

\example{Head with an explicit nonempty precondition}{\Uone{}+\Dep{}}
\par\noindent\begin{minipage}{\linewidth}
\begin{lstlisting}
head(A, n in int & >=1): func(xs: Vec(A, n)) -> A: xs[0]
first: head([4, 5])         // 4
bad: head([])               // no admissible instance
\end{lstlisting}
\end{minipage}\par

\example{Appending adds lengths}{\Uone{}+\Dep{}}
\par\noindent\begin{minipage}{\linewidth}
\begin{lstlisting}
append(A, n in Nat, m in Nat): func(
    xs: Vec(A, n), ys: Vec(A, m),
) -> Vec(A, n + m): [for x in xs {x}, for y in ys {y}]

v: append([1, 2], [3])      // [1, 2, 3]
\end{lstlisting}
\end{minipage}\par
The proof combines the cardinalities of two concatenated finite comprehensions.

\example{Zip rejects unequal lengths}{\Uone{}+\Dep{}}
\par\noindent\begin{minipage}{\linewidth}
\begin{lstlisting}
zip(A, B, n in Nat): func(xs: Vec(A, n), ys: Vec(B, n)) ->
    Vec([A, B], n): [for i, x in xs {[x, ys[i]]}]

z: zip([1, 2], ["a", "b"]) // [[1, "a"], [2, "b"]]
bad: zip([1], ["a", "b"])   // inconsistent length witness
\end{lstlisting}
\end{minipage}\par
The same existential instance of $n$ must satisfy both input constraints at a
call. Separate instances per argument would lose this invariant.

\example{Split at a bounded position}{\Uone{}+\Dep{}}
\par\noindent\begin{minipage}{\linewidth}
\begin{lstlisting}
split(A, n in Nat, k in int & >=0 & <=n): func(
    xs: Vec(A, n), cut: k,
) -> {left: Vec(A, k), right: Vec(A, n - k)}: {
    left: xs[:cut]
    right: xs[cut:]
}
\end{lstlisting}
\end{minipage}\par
The runtime parameter \code{cut} reifies the index used by slicing. The endpoints
$k=0$ and $k=n$ are admitted.

\example{Replication reifies its count}{\Uone{}+\Dep{}}
\par\noindent\begin{minipage}{\linewidth}
\begin{lstlisting}
replicate(A, n in Nat): func(x: A, count: n) -> Vec(A, n): [
    for i in list.Range(0, count, 1) {x}
]
r: replicate("x", 3)      // ["x", "x", "x"]
\end{lstlisting}
\end{minipage}\par
An erased type or value index cannot be inspected by a body merely because it
appears in a contract. Here the actual count parameter supplies the needed
runtime information. A checked library bridge can carry its declared allocation
or representation effects; the finite reference range operation is total.

\example{A dependent result with a hidden length}{\Shigher{}+\Dep{}}
\par\noindent\begin{minipage}{\linewidth}
\begin{lstlisting}
filterBounded(A, n in Nat): func(
    p: func(A) -> bool, xs: Vec(A, n),
) -> (exists (k in int & >=0 & <=n) Vec(A, k)):
    [for x in xs if p(x) {x}]
\end{lstlisting}
\end{minipage}\par
The existential witness can depend on the input and predicate. A \Uone{}+\Dep{} surface
can express the same projected result using \code{[...A] & list.MaxItems(n)}.
A witness-bearing version returns a record \code{{length: k, values: Vec(A,k)}}.

\section{Matrix dimensions, including empty matrices}
Retain dimensions as data when empty collections do not determine them:
\par\noindent\begin{minipage}{\linewidth}
\begin{lstlisting}
Matrix(A, r in Nat, c in Nat) = {
    rows: r
    cols: c
    data: Vec(Vec(A, c), r)
}
\end{lstlisting}
\end{minipage}\par
A zero-row matrix still records its column count. This makes transposition and
composition defined at all admitted dimensions.

\example{Transpose with a precise shape}{\Uone{}+\Dep{}}
\par\noindent\begin{minipage}{\linewidth}
\begin{lstlisting}
transpose(A, r in Nat, c in Nat): func(m: Matrix(A, r, c)) ->
    Matrix(A, c, r): {
        rows: m.cols
        cols: m.rows
        data: [
            for j in list.Range(0, m.cols, 1) {
                [for row in m.data {row[j]}]
            }
        ]
    }
\end{lstlisting}
\end{minipage}\par
Every $j$ is below the length $c$ of every row. The outer length is $c$ and the
inner length is $r$, including either zero dimension.

\example{Dot product with matched dimensions}{\Uone{}+\Dep{}; finite fold primitive}
\par\noindent\begin{minipage}{\linewidth}
\begin{lstlisting}
addNumber: func(x: number, y: number) -> number: x + y

dot(n in Nat): func(xs: Vec(number, n), ys: Vec(number, n)) ->
    number: fold(addNumber, 0, [for i, x in xs {x * ys[i]}])
\end{lstlisting}
\end{minipage}\par
Products and sums return \code{number}; arbitrary input refinements are not
assumed closed under arithmetic.

\example{Dimension-safe matrix multiplication}{\Uone{}+\Dep{}}
\par\noindent\begin{minipage}{\linewidth}
\begin{lstlisting}
matmul(r in Nat, k in Nat, c in Nat): func(
    x: Matrix(number, r, k), y: Matrix(number, k, c),
) -> Matrix(number, r, c): {
    rows: x.rows
    cols: y.cols
    data: [
        for row in x.data {
            [for j in list.Range(0, y.cols, 1) {
                dot(row, [for yrow in y.data {yrow[j]}])
            }]
        }
    ]
\end{lstlisting}
\end{minipage}\par
A single instance of $k$ links the two inputs and justifies each dot product.
The output carries the surviving dimensions $r,c$.

\section{Protocol descriptions and refinement predicates}
\example{A request-indexed response}{\Uone{}+\Dep{} with a checked service}
\par\noindent\begin{minipage}{\linewidth}
\begin{lstlisting}
let Request = {kind: "lookup", key: string} |
              {kind: "size", items: [..._]}
Reply(q in Request) = {
    if q.kind == "lookup" {result: string}
    if q.kind == "size"   {result: int & >=0}
}
request: extern func(q: Request) -> Reply(q) !service
\end{lstlisting}
\end{minipage}\par
The result predicate is evaluated in the request's environment. Replacing it by
an independent union of responses would admit mismatched request/response pairs.

\example{A relation between input and output records}{\Uone{}+\Dep{}}
\par\noindent\begin{minipage}{\linewidth}
\begin{lstlisting}
annotate(A): func(x: A, tag: string) -> {
    value: x
    label: tag
}: {value: x, label: tag}
\end{lstlisting}
\end{minipage}\par
The field labels differ from the parameter names. The result predicate states
equality with the corresponding argument cells and retains the input/output
correlation through subsequent record refinement.

\example{Clamping to an argument-dependent interval}{\Uone{}+\Dep{}}
\par\noindent\begin{minipage}{\linewidth}
\begin{lstlisting}
clamp: func(lo: int, hi: int & >=lo, x: int) -> (
    int & >=lo & <=hi
): {
    if x < lo {out: lo}
    if x > hi {out: hi}
    if x >= lo && x <= hi {out: x}
}.out

bounded: clamp(0, 10, 14) // 10
\end{lstlisting}
\end{minipage}\par
The three branches partition the input domain. The dependency $hi\ge lo$
justifies each branch's interval postcondition. The result is established by
linear arithmetic and branch refinement.

\section{Proof fragments and erasure}
Dependent specifications can be verified by a hierarchy of procedures:
structural length reasoning, quantifier-free linear integer arithmetic,
finite disjunction splitting, and explicit proof certificates for stronger
predicates. Refinement-type inference provides a precedent for combining type
structure with restricted automated implication checking \cite{liquid}.

The logical reference model admits every well-formed predicate. An executable
profile grants a contract only after a sound derivation. A residual obligation
retains its variables and scope. Runtime checks produce declared errors and do
not turn a conditional check into a successful-call theorem.

The erasure property states that changing a proved annotation or choosing a
type-instantiation derivation does not change the underlying closure and its
runtime call packet. Value indices used computationally are ordinary arguments
or fields, as in replication and matrices. This separates specification
parameters from executable data without changing their shared set semantics.

\chapter{Existentials and modular reuse}
\label{ch:modules}
\section{A shared hidden representation}
An existential record describes several fields using one representation witness:
\par\noindent\begin{minipage}{\linewidth}
\begin{lstlisting}
#Counter: exists State {
    zero: State
    next: func(State) -> State
    read: func(State) -> (int & >=0)
}
\end{lstlisting}
\end{minipage}\par
Its meaning is $\bigcup_{S\in\U_n}\sem{\operatorname{Counter}(S)}$.
One $S$ must make all three fields satisfy the interface. Replacing it with
separate existential witnesses for the three fields loses the relationships
that make composition type-correct.

This is the type-theoretic basis of abstract data types and modules identified
by Mitchell and Plotkin \cite{mitchellplotkin}. In CUE, an additional distinction
matters: existential projection is a logical operation, while \emph{information
hiding} is a restriction on the observations available at an interface.

\section{Sealing and opening}
A \emph{seal} turns an implementation of an existential interface into an opaque
public view. For an unbounded representation sort $S$, a private representation
$A$ is exposed through a fresh abstract copy
\[
 \widehat A_\kappa=\{(\kappa,a)\mid a\in A\},
\]
where $\kappa$ is a seal identity. Exported operations wrap and unwrap that copy
through authorized adapters. The public existential witness is $\widehat A_\kappa$.
Thus the set interpretation of the interface remains an existential union.

Public operations are exposed through opaque capability handles associated
with the seal and the interface's declared export/alias graph. Their private
code origins and captures are not observable through that view. Copying a sealed
module preserves its seal and exported aliasing; explicitly constructing a fresh
sealed module introduces a fresh seal.

Sealing and opening have the forms
\par\noindent\begin{minipage}{\linewidth}
\begin{lstlisting}
seal Interface with (HiddenType = Representation) {
    // implementation fields, checked with the private representation
}

open moduleValue as (AbstractType, View) {
    // AbstractType is rigid; View has the opened interface.
}
\end{lstlisting}
\end{minipage}\par
Opening introduces a local rigid type and preserves sharing across the view's
field projections. The type cannot escape except under a corresponding
existential package. \Sone{} uses this operation at module boundaries; \Shigher{} admits it
inside functions over first-class packages.

A transparent bounded existential still exposes every operation justified by its
bound. A fully opaque initial sealing profile uses unbounded representation
sorts and explicit exported operations. A bound such as \code{State: number}
would expose numeric structure and entails a stronger observation-preservation
condition on any representation change.

\section{Two implementations and one client}
\example{An integer representation}{\Sone{}+\Abs{}}
\par\noindent\begin{minipage}{\linewidth}
\begin{lstlisting}
counterV1: seal #Counter with (State = int & >=0) {
    zero: 0
    next: func(x: int & >=0) -> (int & >=0): x + 1
    read: func(x: int & >=0) -> (int & >=0): x
}
\end{lstlisting}
\end{minipage}\par

\example{A record representation}{\Sone{}+\Abs{}}
\par\noindent\begin{minipage}{\linewidth}
\begin{lstlisting}
let CounterRep = close({count: int & >=0})
counterV2: seal #Counter with (State = CounterRep) {
    zero: {count: 0}
    next: func(s: CounterRep) -> CounterRep:
        {count: s.count + 1}
    read: func(s: CounterRep) -> (int & >=0): s.count
}
\end{lstlisting}
\end{minipage}\par
The abstract carrier in each public view has no accessible \code{count} field
or numeric representation. Its structure is the interface's structure.

\example{A client checked once}{\Sone{}+\Abs{} with a rank-1 utility}
\par\noindent\begin{minipage}{\linewidth}
\begin{lstlisting}
useCounter(S): func(c: {
    zero: S
    next: func(S) -> S
    read: func(S) -> (int & >=0)
}) -> (int & >=0): c.read(c.next(c.next(c.zero)))

v1: (open counterV1 as (S, C) {out: useCounter[S](C)}).out
v2: (open counterV2 as (S, C) {out: useCounter[S](C)}).out
// Both observations are 2.
\end{lstlisting}
\end{minipage}\par
The client uses only the abstract operations. Its generic declaration is
rank-1. The hidden type is opened at the module boundary and then supplied as a
rigid instance.

\example{A module law strengthened by universals}{\Sone{}+\Dep{}+\Abs{}}
\par\noindent\begin{minipage}{\linewidth}
\begin{lstlisting}
#CounterLawful: exists State {
    zero: State
    next: func(State) -> State
    read: func(State) -> (int & >=0)
    _zeroLaw: true & (read(zero) == 0)
    forall (s in State) {
        _stepLaw: true & (read(next(s)) == read(s) + 1)
    }
}
\end{lstlisting}
\end{minipage}\par
The law is a proof obligation on the module subject, not an instruction to
enumerate all states. Both implementations discharge it by substitution into
their bodies. The existential representation scope encloses the universal
state-law scope.

\section{Unification is an observation}
Opaque values can participate in CUE unification. In the basic sealed model,
\[
 (\kappa,a)\mathbin{\&}(\kappa,b)
 =\begin{cases}(\kappa,a),&a=b,\\\bot,&a\ne b.\end{cases}
\]
Consequently, a representation-independence theorem must preserve equality and
conflict, not merely the outputs of named methods.

\begin{definition}[Equality-respecting representation simulation]
For two representations $A_1,A_2$, a simulation is a bijection
$r:A_1\to A_2$ on the entire carriers admitted by the public interface,
together with corresponding exported capability handles, such that constants
correspond, every public operation commutes with $r$, and public data results
and effects agree. The carriers include the abstract states over which a client
may quantify, as well as those reached by executing operations.
\end{definition}
For the counters, $r(n)=\{\mathit{count}:n\}$ is such a bijection. It preserves
zero, successor, reading, and equality. The abstract public seal identities are
related consistently across the two executions.

Extend $r$ componentwise to data constructors and to sets of completions by
direct image. A bijection preserves intersections, unions, emptiness, and
singletons. Its transport of abstract type assignments induces bijections on
the corresponding admissible predicate families.

\begin{theorem}[Representation-independent replacement]\label{thm:abstraction}
Assume an equality-respecting simulation, scope-preserving sealing and opening,
and admissible type families and public operational inhabitants closed under
the induced transport. Let the client's observations be confined to the declared
public interface, ordinary data, unification, and the admitted effects. Replacing one sealed
implementation by the related implementation preserves every client data
observation and its success/conflict outcome.
\end{theorem}
\begin{proof}
Interpret the client in the two related public environments. Constants and
public operations preserve the relation by hypothesis. Reindexing and data
constructors preserve it componentwise. Direct image along a bijection commutes
with unions and intersections. Quantifiers range over corresponding assignments,
so their unions and intersections commute with transport. Application preserves
the relation by the operation simulation. Sealing preserves the public alias
structure and hides implementation descriptors. Induction over the elaborated
client derivation therefore relates its resulting completion predicates. At
ordinary exported data, the relation is equality; emptiness also agrees.
\end{proof}
The abstraction discipline is related to the classical logical-relations proof
of representation independence \cite{mitchell86,reynolds83}. The explicit
equality condition is required by CUE's observable meet operation.

A many-to-one abstraction relation requires additional treatment. If multiple
private cache states represent the same logical state, their raw equality can
be observable through unification. An interface may instead expose a certified
quotient carrier and an equality congruence. Its sealing adapter then presents
logical equivalence classes, and the theorem applies to the resulting public
carrier. The implementation must realize that equality soundly.

\section{The scope of the evolution guarantee}
The theorem yields a checkable class of representation changes for which all
admissible clients continue to work. Its hypotheses include behavior and
observable equality. A signature that merely types \code{next} and \code{read}
does not assert that \code{read(next(s)) = read(s)+1}; that equation is a separate
interface law or simulation obligation.

For one-way evolution, an observation-preserving simulation can relate the old
interface to a designated old-compatible view of the new implementation. New
operations and states may be exposed through an extended view. If additional
states are exposed through the old abstract type itself, preservation also
requires every old universally quantified obligation to hold on those states;
agreement merely on previously executed states is insufficient. Effect
guarantees and public serialization are part of the same simulation.

Here a \emph{forward-compatible client} is a client whose old interface contract
continues to hold when a later, related implementation is substituted. This is
the usual backward-compatible replacement direction for the library. The
quantified interface and simulation state the precise class of permitted future
changes.

Ordinary exposed structural schemas already support structural refinement and
interface checks. Sealed existential witnesses add a generic
\emph{representation-independent} replacement guarantee. A private field alone
does not abstract the representation of a public value whose numeric, record,
or list structure remains observable.

\section{First-class packages and existential placement}
\example{Heterogeneous values with their own operations}{\Shigher{}+\Abs{}}
\par\noindent\begin{minipage}{\linewidth}
\begin{lstlisting}
#Showable: exists A {
    value: A
    show: func(A) -> string
}
p: seal #Showable with (A = int) {
    value: 7
    show: func(x: int) -> string: "\(x)"
}
q: seal #Showable with (A = string) {
    value: "hello"
    show: func(x: string) -> string: x
}
items: [p, q]
\end{lstlisting}
\end{minipage}\par
Each package has its own type witness and seal. A list of packages differs from
one package containing a homogeneous list.

\example{Eliminating a first-class package}{\Shigher{}+\Abs{}}
\par\noindent\begin{minipage}{\linewidth}
\begin{lstlisting}
describe: func(p: #Showable) -> string:
    (open p as (A, P) {out: P.show(P.value)}).out
text: map(describe, items)     // ["7", "hello"]
\end{lstlisting}
\end{minipage}\par
At a sealed interface, projections are mediated by the opened abstract view.
The checker does not independently instantiate the hidden type for \code{show}
and \code{value}.

\example{Sharing two values under one witness}{\Shigher{}+\Abs{}}
\par\noindent\begin{minipage}{\linewidth}
\begin{lstlisting}
#ComparablePair: exists A {
    left: A
    right: A
    equal: func(A, A) -> bool
}
compare: func(p: #ComparablePair) -> bool:
    (open p as (A, P) {out: P.equal(P.left, P.right)}).out
\end{lstlisting}
\end{minipage}\par
Separate existential packages for the two operands would require an explicit
sharing certificate before one comparator could consume both.

\example{An abstract codec interface}{\Sone{}+\Abs{}; checked effects}
\par\noindent\begin{minipage}{\linewidth}
\begin{lstlisting}
#CodecModule: exists Value {
    encode: func(Value) -> bytes
    decode: func(bytes) -> Value !decode
}
\end{lstlisting}
\end{minipage}\par
A round-trip law requires decoding an encoded value to return the corresponding
abstract value. Persisted wire compatibility additionally constrains the chosen
byte representation; hiding an in-memory representation does not silently
version an external protocol.

\Needspace{180pt}
\section{Mixed prefixes and generic modules}
The order of module binders expresses representation dependence:
\par\noindent\begin{minipage}{\linewidth}
\begin{lstlisting}
module(A): {
    exists State
    empty: State
    push: func(A, State) -> State
}
\end{lstlisting}
\end{minipage}\par
The logical order is $\forall A\,\exists State$. It allows the hidden semantic
state type to depend on $A$ while constraining one uniform module subject. A
standard witness is $State=\List(A)$, with empty list and a generic cons operation.
It differs from $\exists State\,\forall A$, which chooses one semantic state type
before the element type varies.

When different \emph{runtime} modules are required for different arguments, an
ordinary module-producing function supplies the module subject. Quantifiers
specify its interface; they do not themselves allocate a record at each type
assignment. This keeps generativity, runtime allocation, and logical dependence
separate and explicit.

\chapter{Checking, execution, and formal obligations}
\label{ch:checking}
\section{Judgments and rigid scopes}
Use a constraint context $\Gamma$ with rigid variables, flexible existential
variables, and a scope level for each variable. The principal judgments are
\[
 \Gamma\vdash e\Rightarrow T,
 \qquad \Gamma\vdash e\Leftarrow T,
 \qquad \Gamma\models T\le U.
\]
They respectively synthesize a description, check an expression against a
description, and establish semantic inclusion. The checker returns a proof,
a refutation with a witness, or a residual obligation.

Universal introduction checks one expression under an arbitrary rigid variable:
\[
 \frac{\Gamma,A:\Type(n)\vdash e\Leftarrow T
       \quad A\notin\FV(\Gamma,e_{\rm run})}
      {\Gamma\vdash e\Leftarrow\forall A\,T}.
\]
The side condition refers to erased runtime code; $A$ may occur in checked
annotations. A bounded introduction adds its bound as an assumption. The same
runtime expression must be valid at every assignment.

Universal elimination substitutes a well-sorted admissible type. Inference may
introduce a fresh flexible type and solve its constraints. Different uses obtain
different instantiation variables; the declared universal remains unchanged.
Existential introduction supplies a witness. Elimination opens one shared witness
in a fresh rigid scope, and the result must respect its escape condition.

These scope rules have implementation precedents in higher-rank bidirectional
checking and quantified logic programming \cite{dunfield,nadathur}. A flexible
witness may depend only on the rigid variables available at its introduction
site. The solver never resolves an outer existential by referring to a later
universal.

\section{Functions and declaration conjunction}
A lambda rule checks the body on arbitrary admitted packets and proves its
successful result predicate. In the nonrecursive core, evaluation is structurally
bounded; primitives carry checked contracts. A structural-recursion profile
adds termination certificates.

Conjunction accumulates predicates on the same subject. In particular,
\[
 \frac{\Gamma\vdash e\Leftarrow T\qquad\Gamma\vdash e\Leftarrow U}
      {\Gamma\vdash e\Leftarrow T\land U}
\]
concerns one expression, not a synthesis of two bodies. Application uses the
correlation-preserving arrow clauses applicable to the actual packet.
Subsumption uses semantic inclusion. Existential calls can propagate constraints
on their argument and result cells without granting an unproved universal
capability to the selected implementation.

\begin{theorem}[Soundness of the reference rules]\label{thm:soundness}
Assume sound primitive contracts and implication procedures, a scope-correct
elaboration, and the termination obligations of the selected execution profile.
If the reference checker establishes $\Gamma\vdash e\Leftarrow T$, then every
admitted environment gives the erased expression an inhabitant of $\sem T$.
For a checked successful arrow, every admitted call packet produces a nonempty
result description contained in its codomain. Checked effects are limited to
the declared outcome set.
\end{theorem}
\begin{proof}
Induct on the checking derivation. Base cases use the primitive hypotheses.
Conjunction and disjunction use set membership. Subsumption uses inclusion.
The lambda case is the universal packet obligation in \cref{eq:arrow}; the
application case instantiates it. Universal introduction ranges over arbitrary
rigid assignments while preserving erased code; elimination specializes the
intersection. Existential rules choose or open a witness while preserving its
scope. Termination and primitive outcome obligations establish the successful
or admitted-effect cases.
\end{proof}
Sound checking permits annotations and residual obligations; it need not be
complete.

\section{Adequate symbolic evaluation}
Let $\operatorname{Comp}(E)$ denote successful complete observations of an
elaborated expression, and let $N(E)$ be the evaluator's retained symbolic graph.
The adequacy target is
\[
 \sem{N(E)}=\operatorname{Comp}(E)
\]
for the selected reference fragment, with preference and failure observations
recorded separately. A conservative abstract interpreter may instead retain an
over-approximation for propagation; its final conformance judgments must still
be justified independently.

Fresh call activations copy parameter cells and preserve closure captures.
Function application does not invent implementations to satisfy a codomain.
A call on an unspecified implementation can remain incomplete. A contradiction
is established only when the retained constraints have no admissible completion,
or a concrete computation produces a specified failure.

Memoization keys include closure identity, bound packets, and the relevant
resolved environment. Residual obligations are retained with their dependency
links. A later refinement invalidates stale approximations or further restricts
them. Cached concrete results can only survive or conflict under the refinement
theorem.

\section{Decision fragments}
A useful implementation combines finite Boolean type structure, a sound
semantic-subtyping procedure, finite sets of polymorphic instances, and a
restricted solver for dependent predicates. Castagna--Laurent--Nguyen give a
sound terminating reconstruction algorithm for a particular first-order
polymorphic set-theoretic calculus \cite{cln2024}. That result motivates \Uone{};
its hypotheses and arrow interpretation remain part of its own calculus.

An implementation reports \emph{established}, \emph{refuted}, or
\emph{unresolved}. Treating an unresolved implication as false would turn later
proof discovery into an apparent repair of bottom. Treating it as true would
lose the universal guarantee. The three outcomes therefore belong to the
checker contract. \Cref{ch:residual} specifies a finite local kernel, persistent
proof obligations, and demand-driven concrete validation for that contract.

\chapter{Refinement-stable normalization and checking}
\label{ch:residual}
\section{Hypotheses, refutation, and observation}
A declaration contributes a hypothesis to a retained constraint store. Accepting
that store for further unification does not certify that it has a completion.
Denotational satisfiability and operational refutability are distinct judgments:
\[
 \sem P\ne\varnothing,
 \qquad P\vdash_{\mathcal K}\bot.
\]
Here $\mathcal K$ is a specified, sound proof kernel. Soundness requires
$P\vdash_{\mathcal K}\bot\Longrightarrow\sem P=\varnothing$; the converse
is not required. An unrefuted empty predicate remains denotationally empty.
Its residual representation records the hypotheses whose incompatibility has
yet to be established.

\example{An empty arithmetic relation retained as hypotheses}{Base fragment}
\par\noindent\begin{minipage}{\linewidth}
\begin{lstlisting}
a: int & b * 5
b: a - 1
\end{lstlisting}
\end{minipage}\par
The equations imply $a=5(a-1)$, hence $a=5/4$, contradicting $a\in\mathbb Z$.
A kernel that evaluates multiplication and subtraction only when their operands
are concrete retains these equations. It need not perform symbolic elimination.
After adding \code{a: m} for any concrete integer $m$, it computes $b=m-1$
and checks $m=5(m-1)$. That concrete instance is refuted. For example,
\code{a: 0} yields \code{b: -1} and the conflict $0=-5$.
Refuting one instance constrains that instance; it is not a proof that an
independently instantiated schema has no other witnesses.

\example{A local bound contradiction}{Base fragment}
\par\noindent\begin{minipage}{\linewidth}
\begin{lstlisting}
x: <0
y: x & >0
\end{lstlisting}
\end{minipage}\par
Copying the constraint on \code{x} into \code{y} exposes the empty interval
$(-\infty,0)\cap(0,\infty)$. A local bound rule refutes it without choosing a
concrete atom. This is as stable under refinement as the arithmetic refutation;
it is simply cheaper to discover.

The distinction between suspended expressions and demanded concrete evaluation
already appears in CUE's treatment of reference cycles \cite{cuecycles}.
For the extension, the implementation contract specifies which local rules run
and when observations demand their residual obligations. The mathematical
meaning of a declaration does not depend on the amount of search performed.

\section{Stable reductions}
Write $P\simeq Q$ when two residual predicates have the same denotation on the
same exposed interface, with their local witnesses projected according to their
scopes. A normalizing step must satisfy
\begin{equation}
 P\longrightarrow Q\quad\Longrightarrow\quad
 P\simeq Q.
 \label{eq:residual-equivalence}
\end{equation}
Every subsequent conjunct $R$ then gives
\[
 P\land R\simeq Q\land R.
\]
This is the extension-stability condition. A step introducing fresh local
variables must preserve their projection and the free-variable interface.
Steps under quantifiers use equivalence in that predicate fiber, followed by the
quantifier; they preserve guards and dependencies.

Two useful special cases are consequence propagation and certified refutation:
\[
 \frac{P\vdash_{\mathcal K}C}{P\longrightarrow P\land C},
 \qquad
 \frac{P\vdash_{\mathcal K}\bot}{P\longrightarrow\bot}.
\]
Neither permits deleting an unresolved premise. A primitive equality such as
\code{z: x + y}, with \code{x: 2} and \code{y: 3}, can be replaced by
\code{z: 5} while retaining the input bindings. A conjectured solution to an
existential constraint cannot be committed unless equivalence, rather than mere
satisfiability of one branch, justifies the commitment.

\example{A concrete result with a residual guard}{Base fragment}
\par\noindent\begin{minipage}{\linewidth}
\begin{lstlisting}
x: int
r: 3 & (x + 1)
\end{lstlisting}
\end{minipage}\par
The result cell may carry the candidate atom \code{3}, but the equality
$3=x+1$ remains active. Adding \code{x: 2} validates it; adding \code{x: 5}
refutes it. Replacing the whole expression by an unconditional \code{3} would
lose a constraint. A backward rule deriving \code{x: 2} is also sound if the
primitive's integer semantics justifies it; the minimal kernel need not include
that rule.

Monotonicity constrains admissible decisions, not the completeness of the search
for them. Sound symbolic arithmetic, for example, could refute the first example
immediately. A terminating implementation instead saturates a declared finite
rule fragment and retains the remaining program. This makes the criterion
precise: every committed decision has persistent evidence; computational effort
is a separate resource policy.

Semantic entailment from an empty predicate is vacuous. It therefore cannot, by
itself, justify materializing an arbitrary concrete result. Materialization
requires an actual primitive-evaluation or constructor derivation, together with
any still-active guards. An inconsistency proof produces bottom, not invented
data.

\section{Exact residuals and approximate propagation}
An implementation state is
\[
 S=(\Sigma,\Phi,K,W).
\]
$\Sigma$ is the binding and sharing graph, including quantifier levels and
subjects. $\Phi$ contains the exact residual predicates and unevaluated
expressions. $K$ is a graph of checked consequences with their supporting
premises. $W$ is a worklist of enabled propagation, evaluation, and checking
tasks. Only $\Sigma$ and $\Phi$ determine the state's denotation; $K$ and $W$
record evidence and scheduling.

An approximate descriptor $d_x$ for a cell $x$ has a concretization
$\gamma(d_x)\subseteq V$ satisfying
\[
 \operatorname{proj}_x\sem S\subseteq\gamma(d_x).
\]
Shapes, primitive kinds, explicit bounds, and finite constructor alternatives
are useful such descriptors. Their role is abstract interpretation
\cite{cousot77}. Intersecting two sound upper approximations is sound, and an
empty upper approximation proves a conflict. A nonempty upper approximation
proves neither existence of a completion nor realizability of a function.

The exact predicates remain in $\Phi$. Replacing them by the approximation alone
would enlarge their denotation. In particular, a function's parameter descriptor
is not substituted for a shared existential type variable. Bounds on that
variable stay correlated with the function and with every use of the variable.

Each certificate records its scope and assumptions. A certificate valid under a
guard $G$ licenses the consequence $G\Rightarrow C$, not the unconditional $C$.
An implication proved from other interface hypotheses is a conditional proof;
those hypotheses remain attached until independently discharged or preserved in
the exported interface.

\section{A finite local kernel}
The default kernel operates on a shared graph and retains disjunctions without
expanding them into a Cartesian product. It provides the following rules.

\begin{longtable}{@{}P{0.25\textwidth}P{0.69\textwidth}@{}}
\toprule
Rule family & Evidence and action\\\midrule
\endfirsthead
\toprule Rule family & Evidence and action\\\midrule
\endhead
Atoms and kinds & Compare concrete atoms exactly; intersect primitive-kind
sets; reject a required cell with incompatible kinds.\\
Bounds & Intersect explicitly available numeric bounds, retaining strictness
and integer membership; detect an empty interval.\\
Records and lists & Propagate constraints to shared fields and known positions.
Respect field presence, closedness, and list tails. An optional impossible field
forces absence rather than rejecting an otherwise valid record.\\
Equality & Maintain equivalence of cells established by the elaboration.
Propagate constants across those classes; preserve local witness scopes.\\
Ground primitives & Evaluate pure operations when their required operands are
concrete; retain a suspended operation otherwise.\\
Boolean structure & Discharge a conjunct with a proof. Remove a disjunct only
with a refutation of that branch. Refute a disjunction only when all its arms
are refuted.\\
Quantified clauses & Combine equal-scope universal clauses pointwise; instantiate
at explicit or demanded admissible arguments; preserve the original clause.\\
Function obligations & Decompose annotated arrows and structural bodies using
sound introduction rules; retain unresolved coverage, outcome, and dependent
predicate obligations.\\\bottomrule
\end{longtable}

The record and list rules operate on the host elaborator's presence constraints,
not merely on displayed labels. They retain its distinctions between open,
closed, required, optional, hidden, and definition fields. Constructor cycles
likewise use the selected execution discipline; equality propagation does not
impose an unrelated finite-tree occurs check on recursive descriptions.

A shape mismatch is refutable when the conflicting shape is required in every
remaining possibility. For a function contract, refuting its result shape also
requires an admitted call. A successful arrow with empty domain imposes no
result obligation. A possibly empty guarded domain is therefore retained until
its relevant emptiness or inhabitation is established.

\section{Checking goals are residual constraints}
Formation, type checking, and evaluation share $\Sigma$, $\Phi$, and $K$.
Formation checks binding, sort, universe level, and profile restrictions.
Checking an implementation generates predicates on that implementation;
normalization processes them using the same kernel as other constraints.
Normalization, materialization, and certification select which obligations are
demanded. They share the monotone store-and-proof discipline of constraint
programming \cite{ccp91}.

Distinguish three questions:
\[
 \sem S\subseteq\sem T,
 \qquad \sem S\cap\sem T=\varnothing,
 \qquad f\in\sem T.
\]
They ask for inclusion, disjointness, and membership of a particular
implementation. Failure to prove the first is not evidence for the second.
An implementation known only to satisfy $S$ may still satisfy a narrower $T$.

\example{A broad approximation is not a contradiction}{\Uone{}}
\par\noindent\begin{minipage}{\linewidth}
\begin{lstlisting}
f: func(x: number) -> number: 1
f: func(number) -> int              // provable from the body

unknown: func(number) -> number
unknown: func(number) -> int        // a nonempty refinement
\end{lstlisting}
\end{minipage}\par
The checker uses the body of \code{f}, rather than treating its declared result
upper bound as an exact behavior. The second pair constrains an unspecified
implementation and leaves its witness unresolved.

For a known lambda and an obligation $f\in A\Rightarrow B$, introduce an
arbitrary rigid packet $p\in A$, form a fresh activation, and check the body.
A certificate establishes termination, permitted outcomes, and the required
nonempty result predicate for every such packet. Structural constructors,
identity, certified primitive applications, and finite comprehensions supply
cheap introduction rules. Unsupported arithmetic facts become residual goals.

For example, the kernel can prove arrow inclusion by the rule
\[
 \frac{\Gamma\vdash A\le C\qquad\Gamma\vdash D\le B}
      {\Gamma\vdash(C\Rightarrow D)\le(A\Rightarrow B)}.
\]
For a lambda it can instead prove the stronger property directly from the body,
even when a previously inferred upper bound gives no useful inclusion.
These are alternative proof procedures for the same target predicate.

Proof-search premises and program constraints occupy different parts of the
state. A sufficient rule's undecided premises remain tasks in $W$; they are not
conjoined into $\Phi$ as additional requirements unless they are equivalent to
the original target or are themselves entailed. Otherwise an incomplete proof
attempt would strengthen the program and discard valid implementations.
For a universal goal, a residual premise obtained under a rigid $A$ likewise
retains its universal guard. Dropping that scope would turn an arbitrary-case
obligation into a satisfiable-instance search.

The target obligation may be used as a hypothesis to constrain calls, but not
as its own evidence of implementation conformance. Proof dependencies must be
well founded. An annotated callback parameter is a legitimate assumption of the
caller's conditional proof; the caller's implementation must preserve the
obligation to validate the actual callback. Mutually assuming two implementations'
annotations is not a discharge of either obligation.

\example{Concrete counterexamples to implementation contracts}{\Uone{}}
\par\noindent\begin{minipage}{\linewidth}
\begin{lstlisting}
wrongResult: func(int) -> string
wrongResult: func(x: int) -> int: x

wrongDomain: func(number) -> string
wrongDomain: func(x: int) -> string: "ok"
\end{lstlisting}
\end{minipage}\par
The input \code{0} refutes the first contract. The input \code{1.5} is admitted
by the second contract and rejected by its fixed implementation protocol.
These are finite, stable witnesses, not failed attempts to prove subtyping.

\section{Universal and existential proof search}
A universal clause is stored as a scoped template. It has two complementary uses:
\emph{propagation} instantiates a hypothesis at a demanded argument;
\emph{certification} checks one implementation under arbitrary rigid binders.
A finite set of successful tests supports propagation at those instances; it
does not certify the universal clause.

\example{Uniform proofs and finite refutations}{\Uone{}}
\par\noindent\begin{minipage}{\linewidth}
\begin{lstlisting}
id(A): func(x: A) -> A: x        // one structural proof
bad(A): func(x: A) -> A: 0      // refuted at A = 1, x = 1
\end{lstlisting}
\end{minipage}\par
For \code{id}, the result is the input cell, so the proof uses the same arbitrary
$A$. For \code{bad}, the singleton instance gives incompatible atoms. A kernel
may search a small, fixed vocabulary of witness types and values. Failure of
that search leaves the obligation unresolved.

\example{A symbolic inconsistency exposed by one universal instance}{\Uone{}}
\par\noindent\begin{minipage}{\linewidth}
\begin{lstlisting}
f(A): func(A) -> A
r: f(3) & string
\end{lstlisting}
\end{minipage}\par
The instance $A=\{3\}$ entails that a successful result is \code{3}. Together
with the required result constraint, this refutes the configuration without
selecting an implementation for \code{f}. The universal hypothesis is retained
throughout; it is not replaced by this one instance.

Existential introduction instead records an unresolved witness and its
constraints. A witness assignment can replace that variable only when the
substitution is an exact solution or when projection justifies the elimination.
A search heuristic that finds one of several possible witnesses retains the
alternatives as a disjunction or a residual existential. It cannot choose one
arbitrarily in an order-independent configuration.

\example{A flexible witness must remain refinable}{\Sone{}}
\par\noindent\begin{minipage}{\linewidth}
\begin{lstlisting}
x: exists (n in 1 | 2) {value: n}
x: {value: 2}
\end{lstlisting}
\end{minipage}\par
The first declaration admits both witnesses. Committing to \code{n = 1} would
introduce an error repaired by choosing a different solution. Retaining the
existential allows the second declaration to remove the first alternative.

Rigid parameters are arbitrary but not concrete. Their operations remain
symbolic except where the kernel proves a uniform result. Their scope levels
prevent an outer existential from being solved by a later universal. Binder
normalization uses the laws of \cref{ch:refinement}, including the empty-range
side condition for universal floating.

\section{Counterexamples and remaining environments}
A failed candidate is not necessarily a failed configuration. Suppose a guarded
universal is $\forall a\in D_\rho\,P(\rho,a)$ and a concrete $d$ has been
identified. An unconditional refutation requires evidence, under the current
remaining environments, that $d\in D_\rho$ and $\neg P(\rho,d)$.
If failure holds only for a subset $H$ of those environments, the sound
consequence is exclusion of $H$. Other environments remain possible.

\example{A refinable argument restricts a call, not a universal domain}{\Uone{}}
\par\noindent\begin{minipage}{\linewidth}
\begin{lstlisting}
x: number
f: func(x) -> string
f: func(n: int) -> string: "ok"

x: int
\end{lstlisting}
\end{minipage}\par
The singleton argument contract on \code{x} excludes noninteger assignments.
A trial with \code{x = 1.5} refutes that assignment, not every assignment of
\code{x}. Refining it to \code{int} cannot repair bottom, because the earlier
joint constraint was never refuted. The same rule applies to an existential
type variable $X$ in $f\in X\Rightarrow B$: a rejected packet removes choices
of $X$ that contain it, rather than substituting the current upper bound for $X$.

This distinction also governs branch pruning. A counterexample to one candidate
implementation or one union arm does not establish that their existential
union is empty. Branch-local certificates keep their guards until all
alternatives are excluded or a remaining branch is forced.

\section{Demanded calls and concrete observations}
A demand identifies an observation and its dependency graph. Requiring a
concrete data result wakes its primitive operations, equalities, field-presence
checks, and result constraints. An operation producing the wrong concrete atom
or shape yields a refutation of that demanded instance.

The evaluator returns one of
\[
 \mathsf{Conflict}(\pi),\qquad
 \mathsf{Value}(v;S'),\qquad
 \mathsf{Blocked}(S').
\]
$\pi$ is a checked refutation. In $\mathsf{Value}(v;S')$, $v$ has a concrete
constructor or evaluation derivation and $S'$ retains its guards and all
unresolved declarations. It is a concrete observation \emph{with a residual
program}, not a certificate of global satisfiability. A data export declared
complete must discharge its demanded validation obligations; otherwise it
reports incompleteness. A universally specified module declared verified must
also discharge its universal conformance obligations.

\example{Local execution and universal verification}{\Uone{}}
\par\noindent\begin{minipage}{\linewidth}
\begin{lstlisting}
externalArithmetic: extern func(int) -> int !bridge
f: func(x: int) -> int: externalArithmetic(x)
r: f(0)
\end{lstlisting}
\end{minipage}\par
A terminating, correctly typed observation of \code{externalArithmetic(0)} may
establish the particular result of \code{r}, with the remaining contract still
in the store. It does not establish correct termination and results at every
integer. The foreign operation is evaluated only at its declared execution
boundary. Certification of the pure arrow requires independent evidence of the
universal property; a checked foreign arrow instead records the permitted
failures as in \cref{ch:functions}.

A complete closure descriptor is concrete as an implementation identity, but
membership in an arbitrary universal type can still require an unresolved
proof. Concrete-data validation and universal implementation certification have
different completeness requirements.

\begin{proposition}[Concrete-data completeness of the ground kernel]
\label{prop:ground-complete}
Consider a finite, ground validation graph whose leaves are decidable primitive
data predicates, whose operations terminate on the supplied concrete arguments,
and whose required structure and labels are fixed. Fair evaluation and exact
local validation either establish every demanded constraint or produce a finite
refutation. No inconsistent assignment in this fragment is reported complete.
\end{proposition}
\begin{proof}
Evaluate each ready primitive with exact arguments and compare its result with
all required constraints. The graph is finite, and the stated evaluation
hypothesis makes every demanded operation decidable. A false leaf or conflicting
result supplies a refutation; if every check succeeds, their finite conjunction
has been validated at the supplied assignment.
\end{proof}
For an already assigned cyclic graph, validation checks its equations against
that assignment without recursively reevaluating the cycle. Ungrounded cycles
and hidden quantified obligations remain residual unless another rule decides
them. The proposition is a completeness guarantee for checking supplied data,
not for finding its witnesses.

\section{Incremental algorithm}
Each declaration contributes scoped constraint nodes and enabled tasks. The
normalizer is a worklist interpreter with checked rule instances:
\par\noindent\begin{minipage}{\linewidth}
\begin{lstlisting}[language={},basicstyle=\ttfamily\fontsize{8.5}{10.4}\selectfont]
refine(state, declarations, demands, resources):
    elaborate_and_conjoin(state, declarations)
    enqueue_new_dependencies(state, demands)
    while an enabled task has an available resource ticket:
        task = choose_task_deterministically(state)
        result = local_kernel(task, state)
        match result:
            equivalent_rewrite(replacement, proof):
                verify(proof); install_preserving_scopes(replacement)
            consequence(fact, proof):
                verify(proof); conjoin_and_record(fact, proof)
            refutation(guard, proof):
                verify(proof); exclude_guarded_case(guard, proof)
            blocked(dependencies):
                retain_and_watch(task, dependencies)
        enqueue_tasks_whose_inputs_gained_information(state)
    return state_with_all_residuals_and_certificates(state)
\end{lstlisting}
\end{minipage}\par
A proof must be checked before its consequence is committed. A refuted
unconditional case becomes bottom; a guarded case is eliminated only within
its guard. Unchanged blocked tasks sleep until a relevant premise is refined.
The worklist does not repeatedly unfold unresolved arithmetic cycles.

Memoization uses stable subject identities, binder identities, explicit
instantiations, and supporting premises. A proved consequence stays valid under
additional conjunction. An unresolved cache entry is merely a suspended task,
not a cached negative answer. A cached concrete atom keeps its validation edges.

Resource tickets bound speculative branch splitting, witness search, and type
instantiation separately from ground evaluation. Tickets are attached to stable
tasks or origins; adding an unrelated declaration does not withdraw earlier
certificates or steal a completed task's proof. An incremental run preserves
all earlier evidence. A fresh run can replay that evidence independently of its
search budget. File-order independence is a property of the retained predicates
and certificates, not a claim that different search budgets produce identical
printed residual graphs.

\section{Preservation and persistence}
\begin{theorem}[Sound incremental normalization]\label{thm:residual-sound}
Assume scope-correct elaboration, sound local certificates, exact ground
primitives, and retention of residual guards. For every finite normalization
sequence $S\longrightarrow^* S'$:
\begin{enumerate}
\item $\sem S=\sem{S'}$ on the exposed interface;
\item every recorded refutation is sound, and remains sound after additional
conjunction in its recorded scope;
\item every concrete observation retains its validation dependencies, so a
later refinement can preserve its atom or refute that guarded observation, but
cannot justify a different atom for the same retained observation;
\item no universal conformance certificate is obtained solely by enumerating
successful instances or by assuming the obligation being certified.
\end{enumerate}
\end{theorem}
\begin{proof}
Induct on the transitions. Equivalent rewrites use
\cref{eq:residual-equivalence}. Adding a certified consequence preserves the
predicate. A guarded refutation removes only an impossible case. Scope-preserving
quantifier steps use reindexing and the adjunction rules. Ground execution uses
the primitive hypothesis and retains the equations connecting its result to its
inputs. Conjoining another predicate preserves every such equivalence and
implication. A concrete atom is a singleton, so a stronger result constraint can
only retain it or be disjoint from it. The certificate checker enforces the
universal introduction and dependency conditions on every discharged goal.
\end{proof}

The theorem is about finite evidence and residual programs. It permits a
satisfiable store, an unrefuted empty store, and a store awaiting an abstract
implementation to remain incomplete. It does not identify any of them with a
successful closed implementation. Defaults, comprehension elaboration, and
external effects retain their specified observation boundaries.

\section{Cost and profile-specific decision fragments}
For the finite propagation core, fix the graph and the available descriptor
vocabulary. Let $n$ be its number of cells and obligation nodes, $m$ its dependency
edges, $h$ the maximum number of strict descriptor refinements per node, and $c$
a bound on one transfer or certificate check. A worklist that rechecks a node
only when an input changes has the bound
\[
 O\bigl(m\,\alpha(n)+(n+m)h\,c\bigr),
\]
with union--find cost $\alpha$ for equality classes. This bound counts graph
operations; exact arithmetic contributes its bit cost to $c$. Finite bound
endpoints and shape flags give a finite $h$. Ground calls and newly requested
instances extend the graph and are charged separately. Unrestricted interval
strengthening or Boolean distribution does not satisfy the finite-vocabulary
hypothesis.

A small nonbranching kernel is therefore cheap, even though arbitrary semantic
inclusion is not its decision problem. Optional stronger procedures may return
certificates, refine residuals, or exhaust their budget. Exhaustion produces a
residual obligation, never a contradiction or a conformance certificate.

\begin{longtable}{@{}P{0.15\textwidth}P{0.39\textwidth}P{0.40\textwidth}@{}}
\toprule
Profile & Default checking organization & Exact fragment and residual boundary\\\midrule
\endfirsthead
\toprule Profile & Default checking organization & Exact fragment and residual boundary\\\midrule
\endhead
\Uone{} & Check an outer universal once with rigid type variables; create
flexible instances at uses; propagate quantifier-free clause constraints. &
Annotated, acyclic structural terms over a finite base-type algebra admit a
terminating structural checker. Refinement predicates and unresolved semantic
inclusions remain goals.\\
\Sone{} & The same engine with existential module witnesses, scope levels, and
escape checks at fixed boundaries. & Explicit witnesses and openings reduce to
checks in the selected monomorphic fragment. Arbitrary witness synthesis and
module laws remain residual.\\
\Shigher{} & Bidirectional checking with explicit introduction and elimination
at nested quantified positions. & A fixed predicative annotated fragment can
have decidable checking; arbitrary quantified Boolean subtyping is not inferred
from that fact.\\
\Dep{} & Generate value-indexed arithmetic, equality, and shape goals in the
same store. & Ground validation is exact. A selected decidable arithmetic theory
can discharge more goals; unsupported predicates stay symbolic.\\\bottomrule
\end{longtable}

Decidability belongs to a precisely delimited judgment and predicate language,
not to rank alone. Rank-1 set-theoretic reconstruction and annotated higher-rank
bidirectional checking both have terminating algorithms for particular calculi
\cite{cln2024,dunfield}. Their restrictions delimit useful components, rather
than deciding all predicates expressible in the present language.



\chapter{Feature correspondence and related designs}
\label{ch:comparison}
\section{The comparison baseline}
The July 2026 function proposal describes opaque function values, packet binding,
defaults, open signatures, and partial applications. The September theory
addendum analyzes a covariantly tightened call-constraint interpretation and
pointwise conjunction of bodies. It separately proposes coverage obligations
and compatibility checking \cite{cueproposal,cueaddendum}.

The present calculus expresses call restrictions existentially and reusable
capabilities universally. Its set lattice therefore includes both relations
on actual call cells and predicates on implementations. Their different
variance follows from their quantifiers.

\section{Feature coverage}
\small
\begin{longtable}{@{}P{.30\textwidth}P{.64\textwidth}@{}}
\toprule
Feature & Interpretation or construction in quantified CUE\\\midrule
\endfirsthead
\toprule Feature & Interpretation or construction in quantified CUE\\\midrule
\endhead
Legacy records, lists, unions, bounds & Unchanged data predicates and legacy elaboration; \cref{thm:conservative}.\\
Definitions, closedness, optional and required fields & Retained field-presence and label constraints; quantifiers do not add fields.\\
References, embedding, comprehensions & Existing copying plus capture-avoiding binder and witness rebinding.\\
Literal bodies and inferred result constraints & Operational closures with a checked declared or inferred result predicate.\\
Positional and labeled arguments & Complete call packets and a fixed binder.\\
Required, optional, hidden, and anonymous parameters & Presence alternatives, package-qualified labels, and positional slots.\\
Body-only aliases & Lexical aliases with no public packet label.\\
Parameter order and uniqueness & Protocol well-formedness before value unification.\\
Parameter attributes & Retained metadata, interpreted only by the relevant declared extension.\\
Omission defaults & Implementation-owned packet completion; supplied arguments bypass them.\\
Caller disjunction defaults & Existing preference rules on the supplied value.\\
Default contributed by another declaration & Explicit defaulted adapter or a refinable description of a yet-unselected implementation protocol.\\
Label contributed to an existing closure & Explicit label adapter; fresh builtins may receive their generated binder at construction.\\
Function types and intersections & Sets of operational inhabitants satisfying all universal arrow clauses.\\
Open signatures & Existential protocol rows retained until the remaining required packet shape is known.\\
Higher-order parameters and results & Ordinary arrows; universals specify reuse across types.\\
Partial application & A new closure with normalized original-slot bindings and a residual contract.\\
Constraint on a partial value & An ordinary constraint on that residual closure.\\
Distinct bodies combined pointwise & An explicit body that calls both and unifies their result descriptions.\\
Self-unification and captures & Stable origins and structurally constrained captured environments.\\
Validators & Predicates on an existential subject, definable with existing structs.\\
List-based variadic use & One list-valued packet component, with a generic element type.\\
Nonrecursive execution and cycles & The legacy execution profile; structural recursion is a separate opt-in policy.\\
Builtins and foreign calls & Typed descriptors with an explicit admitted effect set and the existing bridge.\\
Compatibility checking & Packet-domain inclusion, output/effect refinement, and module simulations where abstraction applies.\\
Export, formatting, and round trips & Source export retains binder scopes, closure environments, and bound slots; data export retains its existing supported values.\\
Memoization and dependency tooling & Scope-preserving graph traversal and dependency-aware caching.\\
\bottomrule
\end{longtable}
\normalsize

Experimental operations that transform a closure elaborate to explicit
transformations. Contract intersection asserts properties of the original
closure.

\section{Three explicit translations}
\example{Supplying a label}{\Uone{}}
\par\noindent\begin{minipage}{\linewidth}
\begin{lstlisting}
raw: func(_~n: int) -> int: n + 1
named: func(x: int) -> int: raw(x)
a: named(x: 2)              // 3
\end{lstlisting}
\end{minipage}\par
The adapter's packet language admits \code{x}. Raw positional behavior remains
unchanged.

\example{Restricting public calls across files}{\Uone{}}
\par\noindent\begin{minipage}{\linewidth}
\begin{lstlisting}
raw: func(x: number) -> string: "\(x)"
call: {arg: number, result: raw(arg)}
call: {arg: int, result: string}

ok:  (call & {arg: 7}).result   // "7"
bad: (call & {arg: 1.5}).result // _|_
\end{lstlisting}
\end{minipage}\par
This is ordinary CUE instantiation and projection. The argument restriction is
an existential predicate, while \code{raw}'s capability remains universal.

\example{Keeping late contract failures explicit}{\Uone{} with checked effects}
\par\noindent\begin{minipage}{\linewidth}
\begin{lstlisting}
checked: func(x: int) -> string !constraint:
    checkedResult(string, legacyBody(x))
\end{lstlisting}
\end{minipage}\par
Here \code{checkedResult(T, e)} is a reference adapter operation: evaluate $e$,
validate its result against $T$, and return either a successful result or the
named constraint-failure outcome. The adapter's contract describes late failures
honestly. A successful arrow is granted only after their absence is established.

\begin{proposition}[Contract intersection]
Under the universal arrow interpretation, conjoining a contract with a concrete
closure cannot change its call protocol or result. It preserves the singleton
closure or makes its intersection empty.
\end{proposition}
\begin{proof}
A subset of the singleton $\{f\}$ is either $\{f\}$ or empty.
\end{proof}
Consequently, constructing a label/default adapter and asserting a contract are
different operations. Keeping that distinction also preserves associativity.
A full experimental compatibility layer can implement its original invocation
relation as a checked adapter, while translating its public guarantees into the
appropriate quantified predicates.

\section{Modular guarantees}
Universal constraints express a statement the call-restriction interpretation
does not express by itself: one selected implementation handles every admitted
instance. This supports checking \code{map}, polymorphic callbacks, and bounded
refinement-preserving utilities without seeing their eventual calls. The
intersection algebra preserves input/output correlations. Dependent universals
state relationships across argument and result values. Existential seals support
representation-independent replacement under an explicit theorem.

Untyped callbacks and struct encodings can express many of the same
computations. Quantified descriptions additionally preserve their relationships
as constraints.

\section{CDuce and semantic subtyping}
Frisch, Castagna, and Benzaken develop semantic subtyping with function and
Boolean types; Castagna and Xu extend its foundation to polymorphic types and
convex models \cite{semantic,castagnaxu}. CDuce's polymorphic calculi combine
instantiation with intersections and preserve input/output correlations
\cite{poly1,poly2}. These supply close precedents for the arrow and instantiation
layers here.

The current proposal adds explicit quantifiers to the general CUE predicate
language. Its total-success and effect-qualified arrows refer to an operational
inhabitant model. A subtyping algorithm proved for a different arrow model must
be revalidated for the selected reference fragment. Convexity of a semantic
subtyping model and relational parametricity are separate properties.

\section{Elixir and checked dynamic execution}
Elixir's set-theoretic work uses unions, intersections, refinements, and guarded
function behavior. The strong-arrow account pays particular attention to
behavior outside the input domain and interaction with dynamic values
\cite{elixirstrong}. It is relevant to CUE's bridge and legacy-expression
boundaries: failure behavior deserves an explicit interpretation alongside the
successful result type.

This proposal's universals are general constraint binders. Its checked effect
arrows are specified by \cref{ch:functions}, rather than identified with Elixir's
strong-arrow construction. Both approaches use set-theoretic descriptions to
retain more precise information than an uncorrelated union of inputs and outputs.

\section{Hyperdoctrines and compact closure}
Consider objects that are finite interfaces and morphisms that are definable
relations $R\subseteq X\times Y$. Composition is
\[
 (S\circ R)(x,z)\equiv\exists y.\;R(x,y)\land S(y,z).
\]
Juxtaposition of interfaces is the tensor. Equality gives the unit/counit
relations making every object self-dual; the snake identities follow by
existentially eliminating equality witnesses. This is the compact structure of
a relational category.

Adding universal formulas enriches the class of definable relations while
retaining these operations, provided the fragment is closed under the required
composition and equality relations. Conjunction within one predicate fiber is
distinct from the tensor of interfaces. The total executable functions form a
separate subcollection and are not asserted to be compact closed.

Lawvere's adjunctions and the bifibrational account of Melli\`es and Zeilberger
organize these refinement predicates over the relational interface category
\cite{lawvere,mellies}.

\chapter{Requirements and implementation outcome}
\label{ch:requirements}
The addendum's section 9 lists six requirements. The last is its restriction on
expressive strength, referring to section 7. The first five are evaluated below
against the defined reference semantics \cite{cueaddendum}.

\section{Directionality}
A reusable capability is a universal predicate on the implementation. A
particular call is a relation on existential argument and result cells. Neither
judgment guesses whether a field is provided or consumed. A higher-order
parameter introduces its required capability in the caller's type. Module
replacement uses its explicitly declared interface and observation relation.

\section{Monotone refinement}
Conjunction refines the joint predicate. Both quantifiers preserve that
refinement, by \cref{thm:monotone}. Dependency-preserving copying and retention
of bound variables prevent accidental weakening of universal obligations.
An established contradiction remains empty, and an atomic concrete result can
only persist or conflict. \Cref{ch:residual} implements this property by retaining
scoped certificates, residual hypotheses, and validation dependencies. Defaults and comprehension execution retain the
separate observation treatment already identified in the reference model.

\section{Cross-file input and output invariants}
Existing argument/result records express existential call restrictions.
Universally quantified declarations express reusable behavior and dependent
relations. Both can be conjoined across files. \cref{ch:programs,ch:dependent}
give concrete examples of each form, including length preservation and
request-indexed responses.

\section{Callable self-unification}
Function inhabitants have stable semantic origins and captured environments.
Their equality is structural on these descriptors, rather than pointer-based.
Conjunction is idempotent on their singleton predicates, including closures
reached through copying or unrelated record refinement. Partial applications
use normalized original-slot bindings. Sealed operation handles preserve the
interface's public alias structure.

\section{Idiomatic foreign interfaces}
Portable logical domains remain in public signatures. Representability and
conversion failures are explicit checked effects. The existing bridge performs
its existing checks and reports those outcomes. Logical kind/domain mismatches
are distinguished from in-domain foreign failures. Native widths remain private
bridge metadata. No additional conversion stage is required by quantification.

\section{Expressive strength and implementation milestones}
The reference semantics admits higher-rank and dependent specifications. An
execution profile may retain the current nonrecursive policy or add certified
structural recursion. \Uone{} is an implementation entry point with explicit generic
capabilities and semantic arrow intersections. \Sone{} adds abstract type witnesses
at module boundaries. \Shigher{} extends the same formation rules to arbitrary
well-sorted positions.

A conforming implementation establishes the following contracts:
\begin{enumerate}
\item Elaboration preserves legacy data observations, lexical binding, copied
witness sharing, and quantifier dependencies.
\item The checker establishes the reference judgments soundly and keeps
unresolved obligations distinct from contradiction.
\item Symbolic evaluation preserves the successful-completion semantics and the
specified observation/effect behavior.
\item The abstraction boundary preserves seal identity, public aliasing, and
the equality-respecting simulation used for module evolution.
\end{enumerate}

\section{Conclusion}
Universal and existential binders are dual projections of the same relation-valued
constraint semantics. Unification remains intersection. Function types arise by
universally constraining an operational application relation; their variance and
polymorphic behavior follow from that definition. Dependent contexts and
existential modules extend the same structure. The resulting language can be
implemented in rank-restricted stages without changing what its quantified
descriptions mean.

\appendix
\chapter{Core judgments and finite regression model}
\label{app:core}
\section{Formation and binding summary}
A reference type expression has the grammar
\[
\begin{aligned}
 T,U ::= {}& b\mid a\mid T\land U\mid T\lor U\mid\neg T
       \mid T\times U\mid\List(T)\mid T\Rightarrow U\\
       &\mid \{\ell_i:T_i\}_{i\in I}
       \mid \forall a:S\,T\mid\exists a:S\,T.
\end{aligned}
\]
Here $b$ is a base predicate; $a$ is a context variable of the appropriate sort;
$S$ is a well-formed value sort or type universe, possibly with a guard.
Complement is relative to the subject sort. A surface profile may retain only
its supported decidable negative predicates. Dependent descriptions extend $T_i$
with references to earlier argument or record cells, using the shared binding
graph. Function effects refine the application predicate of \cref{ch:functions}.

Contexts are telescopes. The free variables of a binder's sort and guard must
already occur in its context. A bound variable is identified by its binder, not
by its spelling. A quantified expression keeps its subject outside the new
binder; an explicit dependent sum keeps the witness as part of the subject.

\section{Sound introduction and elimination principles}
Write $\Gamma;\Delta\vdash e:T$ for a typing judgment in a term environment
$\Gamma$ and a rigid constraint context $\Delta$. The type variable $a$ is fresh
for $\Gamma$ in universal introduction:
\[
\frac{\Gamma;\Delta,a:S\vdash e:T}
     {\Gamma;\Delta\vdash e:\forall a:S\,T}
\qquad
\frac{\Gamma;\Delta\vdash e:\forall a:S\,T\qquad\Delta\vdash s:S}
     {\Gamma;\Delta\vdash e:T[s/a]}.
\]
The first premise checks the same executable expression at an arbitrary
admissible assignment. Bounds appear as assumptions in $\Delta,a:S$.

For transparent existential descriptions, introduction chooses a witness and
elimination checks the continuation at an arbitrary opened witness:
\[
\frac{\Gamma;\Delta\vdash e:T[s/a]\qquad\Delta\vdash s:S}
     {\Gamma;\Delta\vdash e:\exists a:S\,T}.
\]
Elimination preserves the subject and shares the witness across its projections.
Its result may not expose the local witness without an enclosing existential.
Opaque packages additionally use the sealing discipline in \cref{ch:modules}.

Intersection introduction uses the same term:
\[
\frac{\Gamma;\Delta\vdash e:T\qquad\Gamma;\Delta\vdash e:U}
     {\Gamma;\Delta\vdash e:T\land U}.
\]
Subsumption uses semantic inclusion established by a sound derivation. Application
requires coverage of the complete supplied-call packet and the stated result or
effect obligation. A runtime result check is recorded as an effect unless its
success has been established statically.

The proof of the introduction rules is membership in the corresponding
intersection or union. Elimination uses the selected instance or a scoped
existential witness. These principles specify the reference calculus independently
of a particular constraint-solving strategy.

\section{Regression checks}
The accompanying \code{checks/check_model.py} uses only the Python standard
library. It enumerates four data values, their sixteen subsets, and 625 partial
function tables, with a separate rejection outcome. It executes
\textbf{222,594 assertions} and records the result in \code{checks/results.json}.

\begin{tabularx}{\textwidth}{@{}lY@{}}
\toprule
Family & Property or counterexample\\\midrule
Arrows & Domain-union and result-intersection identities; contravariance.\\
Quantified inhabitants & Singleton-rich universal identity; incompatible
monomorphic refinement.\\
Predicate fibers & Both adjunctions; Frobenius reciprocity; inhabited universal
floating; quantifier monotonicity.\\
Scope and bounds & Mixed quantifiers do not commute; empty-range extrusion
requires its side condition; subtype bounds retain guarded ranges.\\
Application & Refinement inclusion holds; meet equality has a counterexample.\\
Projection & Existential projection need not preserve intersection.\\
Abstract values & Bijection transport preserves meet, join, equality, refinement,
and a finite operation simulation.\\\bottomrule
\end{tabularx}

The finite checks instantiate the stated algebraic laws and counterexamples.
The additional \code{checks/check_residual.py} executes \textbf{245,237 assertions}
for the local kernel of \cref{ch:residual}. It checks preservation of finite
completion sets under bounded propagation, validation of grounded arithmetic
cycles, persistence of refutations, retention of guards and existential
alternatives, and finite quantified-function counterexamples. The results are
recorded in \code{checks/residual-results.json}.
Checker, elaboration, and abstraction obligations are specified in
\cref{ch:checking,ch:residual}.

\chapter{Example and feature index}
\label{app:index}
\begin{longtable}{@{}P{0.35\textwidth}P{0.59\textwidth}@{}}
\toprule
Feature & Location and program forms\\\midrule
\endfirsthead
\toprule Feature & Location and program forms\\\midrule
\endhead
Binding syntax & \Cref{ch:surface}: quantified expressions, declaration
parameters, block prefixes, function-local generics, parametric aliases, and
instantiation selection.\\
Pointwise unification & \Cref{ch:refinement}: alpha-renaming, shared subject,
constant constraints, guarded bounded clauses, non-distributivity.\\
Rank-1 programming & \Cref{ch:programs}: identity, bounded pairing, minimum,
map, filter, partition, flat mapping, composition, adapters.\\
Higher-rank programming & \Cref{ch:programs}: a polymorphic callback used at two
types, returned polymorphic closures, quantified utility records.\\
Calling conventions & \Cref{ch:programs}: all parameter modes, defaults,
omission, hidden labels, partial application, open protocols, attributes.\\
Dependent contracts & \Cref{ch:dependent}: exact successor, interval clamp,
length-preserving map, safe indexing, append, zip, split, replication.\\
Dependent structures & \Cref{ch:dependent}: dimension-carrying matrices,
transpose, dot product, matrix multiplication, indexed service responses.\\
Existential reuse & \Cref{ch:modules}: sealed counters, generic clients, module
laws, heterogeneous packages, shared-witness comparison, codecs.\\
Compatibility & \Cref{ch:comparison}: feature-by-feature correspondence,
adapters, old data code, result combination, experimental migration.\\
Semantic guarantees & \Cref{ch:quantifiers,ch:refinement,ch:functions}:
adjunctions, floating, monotonicity, coverage, variance and application.\\
Residual checking & \Cref{ch:residual}: concrete-data validation, scoped
certificates, local propagation, witness search, and resource bounds.\\
Evolution guarantees & \Cref{ch:modules}: equality-respecting representation
simulation, quantified observations, old-compatible views.\\\bottomrule
\end{longtable}

\clearpage
\renewcommand{\bibname}{References}
\begin{thebibliography}{99}
\addcontentsline{toc}{chapter}{References}
\raggedright

\bibitem{cuespec}
The CUE Authors.
\emph{CUE language specification}.
Living specification, accessed 21 September 2026.
\url{https://cuelang.org/docs/reference/spec/}.

\bibitem{cuerefs}
The CUE Authors.
\emph{CUE language specification: References}.
Specification of copying and rebinding referenced expressions.
Accessed 21 September 2026.
\url{https://cuelang.org/docs/reference/spec/#references}.

\bibitem{cuedefaults}
The CUE Authors.
\emph{Specifying a default value for a field}.
Accessed 21 September 2026.
\url{https://cuelang.org/docs/howto/specify-a-default-value-for-a-field/}.

\bibitem{cuelist}
The CUE Authors.
\emph{Package list}.
Standard-library interface documentation, accessed 21 September 2026.
\url{https://pkg.go.dev/cuelang.org/go/pkg/list}.

\bibitem{cueproposal}
Marcel van Lohuizen.
\emph{Proposal: functions in CUE}. 10 July 2026.
CUE proposal 4484, draft at repository revision
\texttt{a949162aedb86661183c46abef47919d16605014}.
\url{https://github.com/cue-lang/proposal/blob/a949162aedb86661183c46abef47919d16605014/designs/language/4484-functions.md}.

\bibitem{cueaddendum}
Marcel van Lohuizen.
\emph{Why CUE Functions Tighten Instead of Widen}.
Theory addendum to proposal 4484, committed 19 September 2026,
repository revision \texttt{a949162aedb86661183c46abef47919d16605014}.
\url{https://github.com/cue-lang/proposal/blob/a949162aedb86661183c46abef47919d16605014/designs/language/4484-theory-addendum.md}.

\bibitem{lawvere}
F. William Lawvere.
Adjointness in foundations.
\emph{Dialectica}, 23(3--4):281--296, 1969.
Reprinted with commentary in \emph{Reprints in Theory and Applications of
Categories}, 16:1--16, 2006.
\url{https://tac.mta.ca/tac/reprints/articles/16/tr16abs.html}.

\bibitem{mellies}
Paul-Andr\'e Melli\`es and Noam Zeilberger.
\emph{A bifibrational reconstruction of Lawvere's presheaf hyperdoctrine}.
2016. arXiv:1601.06098.
\url{https://arxiv.org/abs/1601.06098}.

\bibitem{reynolds84}
John C. Reynolds.
Polymorphism is not set-theoretic.
In \emph{Semantics of Data Types}, LNCS 173, pages 145--156, 1984.
\url{https://doi.org/10.1007/3-540-13346-1_7}.

\bibitem{reynolds83}
John C. Reynolds.
Types, abstraction and parametric polymorphism.
In \emph{Information Processing 83}, pages 513--523. North-Holland, 1983.
\url{https://www.cs.cmu.edu/~crary/819-f09/Reynolds83.pdf}.

\bibitem{semantic}
Alain Frisch, Giuseppe Castagna, and V\'eronique Benzaken.
Semantic subtyping: Dealing set-theoretically with function, union, intersection,
and negation types.
\emph{Journal of the ACM}, 55(4), Article 19, 2008.
\url{https://doi.org/10.1145/1391289.1391293}.

\bibitem{castagnaxu}
Giuseppe Castagna and Zhiwu Xu.
Set-theoretic foundation of parametric polymorphism and subtyping.
In \emph{ICFP}, pages 94--106, 2011.
\url{https://doi.org/10.1145/2034773.2034788}.

\bibitem{poly1}
Giuseppe Castagna, Kim Nguyen, Zhiwu Xu, Hyeonseung Im, Sergue\"i Lenglet,
and Luca Padovani.
Polymorphic functions with set-theoretic types. Part 1: Syntax, semantics,
and evaluation.
In \emph{POPL}, pages 5--18, 2014.
\url{https://doi.org/10.1145/2535838.2535840}.

\bibitem{poly2}
Giuseppe Castagna, Kim Nguyen, Zhiwu Xu, and Pietro Abate.
Polymorphic functions with set-theoretic types. Part 2: Local type inference
and type reconstruction.
In \emph{POPL}, pages 289--302, 2015.
\url{https://doi.org/10.1145/2676726.2676991}.

\bibitem{cducetutorial}
The CDuce Project.
\emph{Polymorphism tutorial}.
Accessed 21 September 2026.
\url{https://www.cduce.org/tutorial_polymorphism.html}.

\bibitem{cln2024}
Giuseppe Castagna, Micka\"el Laurent, and Kim Nguyen.
Polymorphic type inference for dynamic languages.
\emph{Proceedings of the ACM on Programming Languages}, 8(POPL), 2024.
\url{https://doi.org/10.1145/3632882}.
Preprint: \url{https://arxiv.org/abs/2311.10426}.

\bibitem{elixirstrong}
Jos\'e Valim.
\emph{Strong arrows: a new approach to gradual typing}.
20 September 2023.
\url{https://elixir-lang.org/blog/2023/09/20/strong-arrows-gradual-typing/}.

\bibitem{mitchellplotkin}
John C. Mitchell and Gordon D. Plotkin.
Abstract types have existential type.
\emph{ACM Transactions on Programming Languages and Systems},
10(3):470--502, 1988.
\url{https://doi.org/10.1145/44501.45065}.

\bibitem{mitchell86}
John C. Mitchell.
Representation independence and data abstraction.
In \emph{POPL}, pages 263--276, 1986.
\url{https://doi.org/10.1145/512644.512669}.

\bibitem{dunfield}
Jana Dunfield and Neelakantan R. Krishnaswami.
Complete and easy bidirectional typechecking for higher-rank polymorphism.
In \emph{ICFP}, pages 429--442, 2013.
\url{https://doi.org/10.1145/2500365.2500582}.
Preprint: \url{https://arxiv.org/abs/1306.6032}.

\bibitem{nadathur}
Gopalan Nadathur, Bharat Jayaraman, and Keehang Kwon.
\emph{Scoping constructs in logic programming: Implementation problems
and their solution}.
1998. arXiv:cs/9809016.
\url{https://arxiv.org/abs/cs/9809016}.

\bibitem{liquid}
Patrick M. Rondon, Ming Kawaguchi, and Ranjit Jhala.
Liquid types.
In \emph{PLDI}, 2008.
\url{https://goto.ucsd.edu/~rjhala/papers/liquid_types.html}.

\bibitem{ghcrank}
The GHC Developers.
\emph{GHC User's Guide: Arbitrary-rank polymorphism}.
Accessed 21 September 2026.
\url{https://ghc.gitlab.haskell.org/ghc/doc/users_guide/exts/rank_polymorphism.html}.


\bibitem{cuecycles}
The CUE Authors.
\emph{Reference Cycles}.
CUE language tour, accessed 22 September 2026.
\url{https://cuelang.org/docs/tour/references/cycle/}.

\bibitem{cousot77}
Patrick Cousot and Radhia Cousot.
Abstract interpretation: A unified lattice model for static analysis of programs
by construction or approximation of fixpoints.
In \emph{POPL}, pages 238--252, 1977.
\url{https://www.di.ens.fr/~cousot/COUSOTpapers/POPL77.shtml}.

\bibitem{ccp91}
Vijay A. Saraswat, Martin Rinard, and Prakash Panangaden.
The semantic foundations of concurrent constraint programming.
In \emph{POPL}, pages 333--352, 1991.
\url{https://doi.org/10.1145/99583.99627}.

\end{thebibliography}
\end{document}
